\section{高宗时期其余账本事件锚点}
本节补齐高宗时期尚未导航的账本记录，保留前史、文件、决策和后续四类证据节点的区别。
\subsection{1720—1729}
\eventanchor{EQ-1723-YONGZHENG-ACCESSION-PRELUDE}{「雍正帝即位并开始重整康熙晚…}雍正元年（1723年），「雍正帝即位并开始重整康熙晚年的财政、宗室与官僚议题」之前的条件、信息和争议已通过中央与地方材料显现，构成该事件可追踪的前史节点。核心取证：《清世宗实录》雍正元年相关条；宫中朱批奏折、内阁大库题本与《清史稿》相关志传；同时期地方档案、地方志、日记或对方材料应按具体地区补检。该项锚定的是前史条件与信息背景，不把条件直接写成唯一因果。来源保留：官修编年和地方材料的记录密度与立场并不相同；本行不将条件、传闻、呈报或后见解释直接等同为确定因果。
\eventanchor{EQ-1723-YONGZHENG-ACCESSION-DECISION}{围绕「雍正帝即位并开始重整康…}雍正元年（1723年），围绕「雍正帝即位并开始重整康熙晚年的财政、宗室与官僚议题」的命令、调度、照会或呈报形成独立的中央—地方行动文书链。核心取证：《清世宗实录》雍正元年相关条；宫中朱批奏折、内阁大库题本与《清史稿》相关志传。该项锚定的是命令或决策环节，命令抵达和结果仍须另外核验。来源保留：奏折、上谕、部议、军机档或条约可证明行政动作和机构立场，不自动证明所有地区、群体或对手按其意图行动。
\eventanchor{EQ-1723-YONGZHENG-ACCESSION-AFTERMATH}{「雍正帝即位并开始重整康熙晚…}雍正元年（1723年），「雍正帝即位并开始重整康熙晚年的财政、宗室与官僚议题」引发的善后、执行或地方回应构成可由不同材料追踪的后续节点。核心取证：《清世宗实录》雍正元年相关条；宫中朱批奏折、内阁大库题本与《清史稿》相关志传；相关地区的档案、志书、契约、报刊或对方材料。该项锚定的是后续追踪问题，影响范围须按地区与群体分别判断。来源保留：战后、改革后或条约后的记忆、地方记录和官方总结常具有事后选择性；后续影响须按地点、群体和时间分层，不作整体化推断。
\eventanchor{EQ-1723-JI-ZENGYUN-CASE-PRELUDE}{「年羹尧在西北军政中的权力与…}雍正元年（1723年），「年羹尧在西北军政中的权力与朝廷猜忌加剧」之前的条件、信息和争议已通过中央与地方材料显现，构成该事件可追踪的前史节点。核心取证：《清世宗实录》雍正元年相关条；宫中朱批奏折、内阁大库题本与《清史稿》相关志传；同时期地方档案、地方志、日记或对方材料应按具体地区补检。该项锚定的是前史条件与信息背景，不把条件直接写成唯一因果。来源保留：官修编年和地方材料的记录密度与立场并不相同；本行不将条件、传闻、呈报或后见解释直接等同为确定因果。
\eventanchor{EQ-1723-JI-ZENGYUN-CASE-DECISION}{围绕「年羹尧在西北军政中的权…}雍正元年（1723年），围绕「年羹尧在西北军政中的权力与朝廷猜忌加剧」的命令、调度、照会或呈报形成独立的中央—地方行动文书链。核心取证：《清世宗实录》雍正元年相关条；宫中朱批奏折、内阁大库题本与《清史稿》相关志传。该项锚定的是命令或决策环节，命令抵达和结果仍须另外核验。来源保留：奏折、上谕、部议、军机档或条约可证明行政动作和机构立场，不自动证明所有地区、群体或对手按其意图行动。
\eventanchor{EQ-1723-JI-ZENGYUN-CASE-AFTERMATH}{「年羹尧在西北军政中的权力与…}雍正元年（1723年），「年羹尧在西北军政中的权力与朝廷猜忌加剧」引发的善后、执行或地方回应构成可由不同材料追踪的后续节点。核心取证：《清世宗实录》雍正元年相关条；宫中朱批奏折、内阁大库题本与《清史稿》相关志传；相关地区的档案、志书、契约、报刊或对方材料。该项锚定的是后续追踪问题，影响范围须按地区与群体分别判断。来源保留：战后、改革后或条约后的记忆、地方记录和官方总结常具有事后选择性；后续影响须按地点、群体和时间分层，不作整体化推断。
\eventanchor{EQ-1725-NIAN-GENGYAO-PUNISHMENT-PRELUDE}{「年羹尧被削职、赐死，其党羽…}雍正三年（1725年），「年羹尧被削职、赐死，其党羽与相关官员受到清理」之前的条件、信息和争议已通过中央与地方材料显现，构成该事件可追踪的前史节点。核心取证：《清世宗实录》雍正三年相关条；宫中朱批奏折、内阁大库题本与《清史稿》相关志传；同时期地方档案、地方志、日记或对方材料应按具体地区补检。该项锚定的是前史条件与信息背景，不把条件直接写成唯一因果。来源保留：官修编年和地方材料的记录密度与立场并不相同；本行不将条件、传闻、呈报或后见解释直接等同为确定因果。
\eventanchor{EQ-1725-NIAN-GENGYAO-PUNISHMENT-DECISION}{围绕「年羹尧被削职、赐死，其…}雍正三年（1725年），围绕「年羹尧被削职、赐死，其党羽与相关官员受到清理」的命令、调度、照会或呈报形成独立的中央—地方行动文书链。核心取证：《清世宗实录》雍正三年相关条；宫中朱批奏折、内阁大库题本与《清史稿》相关志传。该项锚定的是命令或决策环节，命令抵达和结果仍须另外核验。来源保留：奏折、上谕、部议、军机档或条约可证明行政动作和机构立场，不自动证明所有地区、群体或对手按其意图行动。
\eventanchor{EQ-1725-NIAN-GENGYAO-PUNISHMENT-AFTERMATH}{「年羹尧被削职、赐死，其党羽…}雍正三年（1725年），「年羹尧被削职、赐死，其党羽与相关官员受到清理」引发的善后、执行或地方回应构成可由不同材料追踪的后续节点。核心取证：《清世宗实录》雍正三年相关条；宫中朱批奏折、内阁大库题本与《清史稿》相关志传；相关地区的档案、志书、契约、报刊或对方材料。该项锚定的是后续追踪问题，影响范围须按地区与群体分别判断。来源保留：战后、改革后或条约后的记忆、地方记录和官方总结常具有事后选择性；后续影响须按地点、群体和时间分层，不作整体化推断。
\eventanchor{EQ-1725-TAN-TEBI-CASE-PRELUDE}{「隆科多失势并被幽禁，雍正初…}雍正三年（1725年），「隆科多失势并被幽禁，雍正初年近臣权力结构继续调整」之前的条件、信息和争议已通过中央与地方材料显现，构成该事件可追踪的前史节点。核心取证：《清世宗实录》雍正三年相关条；宫中朱批奏折、内阁大库题本与《清史稿》相关志传；同时期地方档案、地方志、日记或对方材料应按具体地区补检。该项锚定的是前史条件与信息背景，不把条件直接写成唯一因果。来源保留：官修编年和地方材料的记录密度与立场并不相同；本行不将条件、传闻、呈报或后见解释直接等同为确定因果。
\eventanchor{EQ-1725-TAN-TEBI-CASE-DECISION}{围绕「隆科多失势并被幽禁，雍…}雍正三年（1725年），围绕「隆科多失势并被幽禁，雍正初年近臣权力结构继续调整」的命令、调度、照会或呈报形成独立的中央—地方行动文书链。核心取证：《清世宗实录》雍正三年相关条；宫中朱批奏折、内阁大库题本与《清史稿》相关志传。该项锚定的是命令或决策环节，命令抵达和结果仍须另外核验。来源保留：奏折、上谕、部议、军机档或条约可证明行政动作和机构立场，不自动证明所有地区、群体或对手按其意图行动。
\eventanchor{EQ-1725-TAN-TEBI-CASE-AFTERMATH}{「隆科多失势并被幽禁，雍正初…}雍正三年（1725年），「隆科多失势并被幽禁，雍正初年近臣权力结构继续调整」引发的善后、执行或地方回应构成可由不同材料追踪的后续节点。核心取证：《清世宗实录》雍正三年相关条；宫中朱批奏折、内阁大库题本与《清史稿》相关志传；相关地区的档案、志书、契约、报刊或对方材料。该项锚定的是后续追踪问题，影响范围须按地区与群体分别判断。来源保留：战后、改革后或条约后的记忆、地方记录和官方总结常具有事后选择性；后续影响须按地点、群体和时间分层，不作整体化推断。
\eventanchor{EQ-1726-HAOXIAN-GUIGONG-PRELUDE}{「雍正政府推进耗羡归公与养廉…}雍正四年（1726年），「雍正政府推进耗羡归公与养廉银安排，试图重组地方财政和官员俸给」之前的条件、信息和争议已通过中央与地方材料显现，构成该事件可追踪的前史节点。核心取证：《清世宗实录》雍正四年相关条；户部奏销、盐政、河工或海关档案；同时期地方档案、地方志、日记或对方材料应按具体地区补检。该项锚定的是前史条件与信息背景，不把条件直接写成唯一因果。来源保留：官修编年和地方材料的记录密度与立场并不相同；本行不将条件、传闻、呈报或后见解释直接等同为确定因果。
\eventanchor{EQ-1726-HAOXIAN-GUIGONG-DECISION}{围绕「雍正政府推进耗羡归公与…}雍正四年（1726年），围绕「雍正政府推进耗羡归公与养廉银安排，试图重组地方财政和官员俸给」的命令、调度、照会或呈报形成独立的中央—地方行动文书链。核心取证：《清世宗实录》雍正四年相关条；户部奏销、盐政、河工或海关档案。该项锚定的是命令或决策环节，命令抵达和结果仍须另外核验。来源保留：奏折、上谕、部议、军机档或条约可证明行政动作和机构立场，不自动证明所有地区、群体或对手按其意图行动。
\eventanchor{EQ-1726-HAOXIAN-GUIGONG-AFTERMATH}{「雍正政府推进耗羡归公与养廉…}雍正四年（1726年），「雍正政府推进耗羡归公与养廉银安排，试图重组地方财政和官员俸给」引发的善后、执行或地方回应构成可由不同材料追踪的后续节点。核心取证：《清世宗实录》雍正四年相关条；户部奏销、盐政、河工或海关档案；相关地区的档案、志书、契约、报刊或对方材料。该项锚定的是后续追踪问题，影响范围须按地区与群体分别判断。来源保留：战后、改革后或条约后的记忆、地方记录和官方总结常具有事后选择性；后续影响须按地点、群体和时间分层，不作整体化推断。
\eventanchor{EQ-1727-KYAKHTA-TREATY-PRELUDE}{「《恰克图条约》确定清俄北部…}雍正五年（1727年），「《恰克图条约》确定清俄北部边界与贸易、使团安排」之前的条件、信息和争议已通过中央与地方材料显现，构成该事件可追踪的前史节点。核心取证：《清世宗实录》雍正五年相关条；《筹办夷务始末》相关年卷与中外照会、条约文本；同时期地方档案、地方志、日记或对方材料应按具体地区补检。该项锚定的是前史条件与信息背景，不把条件直接写成唯一因果。来源保留：官修编年和地方材料的记录密度与立场并不相同；本行不将条件、传闻、呈报或后见解释直接等同为确定因果。
\eventanchor{EQ-1727-KYAKHTA-TREATY-DECISION}{围绕「《恰克图条约》确定清俄…}雍正五年（1727年），围绕「《恰克图条约》确定清俄北部边界与贸易、使团安排」的命令、调度、照会或呈报形成独立的中央—地方行动文书链。核心取证：《清世宗实录》雍正五年相关条；《筹办夷务始末》相关年卷与中外照会、条约文本。该项锚定的是命令或决策环节，命令抵达和结果仍须另外核验。来源保留：奏折、上谕、部议、军机档或条约可证明行政动作和机构立场，不自动证明所有地区、群体或对手按其意图行动。
\eventanchor{EQ-1727-KYAKHTA-TREATY-AFTERMATH}{「《恰克图条约》确定清俄北部…}雍正五年（1727年），「《恰克图条约》确定清俄北部边界与贸易、使团安排」引发的善后、执行或地方回应构成可由不同材料追踪的后续节点。核心取证：《清世宗实录》雍正五年相关条；《筹办夷务始末》相关年卷与中外照会、条约文本；相关地区的档案、志书、契约、报刊或对方材料。该项锚定的是后续追踪问题，影响范围须按地区与群体分别判断。来源保留：战后、改革后或条约后的记忆、地方记录和官方总结常具有事后选择性；后续影响须按地点、群体和时间分层，不作整体化推断。
\eventanchor{EQ-1727-BANNER-GENEALOGIES-PRELUDE}{「清廷命令编修八旗谱牒，试图…}雍正五年（1727年），「清廷命令编修八旗谱牒，试图界定旗籍、世系与待遇资格」之前的条件、信息和争议已通过中央与地方材料显现，构成该事件可追踪的前史节点。核心取证：《清世宗实录》雍正五年相关条；地方志、督抚奏折及地方档案；同时期地方档案、地方志、日记或对方材料应按具体地区补检。该项锚定的是前史条件与信息背景，不把条件直接写成唯一因果。来源保留：官修编年和地方材料的记录密度与立场并不相同；本行不将条件、传闻、呈报或后见解释直接等同为确定因果。
\eventanchor{EQ-1727-BANNER-GENEALOGIES-DECISION}{围绕「清廷命令编修八旗谱牒，…}雍正五年（1727年），围绕「清廷命令编修八旗谱牒，试图界定旗籍、世系与待遇资格」的命令、调度、照会或呈报形成独立的中央—地方行动文书链。核心取证：《清世宗实录》雍正五年相关条；地方志、督抚奏折及地方档案。该项锚定的是命令或决策环节，命令抵达和结果仍须另外核验。来源保留：奏折、上谕、部议、军机档或条约可证明行政动作和机构立场，不自动证明所有地区、群体或对手按其意图行动。
\eventanchor{EQ-1727-BANNER-GENEALOGIES-AFTERMATH}{「清廷命令编修八旗谱牒，试图…}雍正五年（1727年），「清廷命令编修八旗谱牒，试图界定旗籍、世系与待遇资格」引发的善后、执行或地方回应构成可由不同材料追踪的后续节点。核心取证：《清世宗实录》雍正五年相关条；地方志、督抚奏折及地方档案；相关地区的档案、志书、契约、报刊或对方材料。该项锚定的是后续追踪问题，影响范围须按地区与群体分别判断。来源保留：战后、改革后或条约后的记忆、地方记录和官方总结常具有事后选择性；后续影响须按地点、群体和时间分层，不作整体化推断。
\eventanchor{EQ-1728-Dzungar-EXPEDITION-PRELUDE}{「清军在西北对准噶尔展开远征…}雍正六年（1728年），「清军在西北对准噶尔展开远征，后勤与高原草原路线成为军事难题」之前的条件、信息和争议已通过中央与地方材料显现，构成该事件可追踪的前史节点。核心取证：《清世宗实录》雍正六年相关条；军机处档、战事方略及地方奏报；同时期地方档案、地方志、日记或对方材料应按具体地区补检。该项锚定的是前史条件与信息背景，不把条件直接写成唯一因果。来源保留：官修编年和地方材料的记录密度与立场并不相同；本行不将条件、传闻、呈报或后见解释直接等同为确定因果。
\eventanchor{EQ-1728-Dzungar-EXPEDITION-DECISION}{围绕「清军在西北对准噶尔展开…}雍正六年（1728年），围绕「清军在西北对准噶尔展开远征，后勤与高原草原路线成为军事难题」的命令、调度、照会或呈报形成独立的中央—地方行动文书链。核心取证：《清世宗实录》雍正六年相关条；军机处档、战事方略及地方奏报。该项锚定的是命令或决策环节，命令抵达和结果仍须另外核验。来源保留：奏折、上谕、部议、军机档或条约可证明行政动作和机构立场，不自动证明所有地区、群体或对手按其意图行动。
\eventanchor{EQ-1728-Dzungar-EXPEDITION-AFTERMATH}{「清军在西北对准噶尔展开远征…}雍正六年（1728年），「清军在西北对准噶尔展开远征，后勤与高原草原路线成为军事难题」引发的善后、执行或地方回应构成可由不同材料追踪的后续节点。核心取证：《清世宗实录》雍正六年相关条；军机处档、战事方略及地方奏报；相关地区的档案、志书、契约、报刊或对方材料。该项锚定的是后续追踪问题，影响范围须按地区与群体分别判断。来源保留：战后、改革后或条约后的记忆、地方记录和官方总结常具有事后选择性；后续影响须按地点、群体和时间分层，不作整体化推断。
\eventanchor{EQ-1729-JUNJICHU-FOUNDING-PRELUDE}{「雍正帝设军机房，后来发展为…}雍正七年（1729年），「雍正帝设军机房，后来发展为军机处，处理紧急军政文书」之前的条件、信息和争议已通过中央与地方材料显现，构成该事件可追踪的前史节点。核心取证：《清世宗实录》雍正七年相关条；宫中朱批奏折、内阁大库题本与《清史稿》相关志传；同时期地方档案、地方志、日记或对方材料应按具体地区补检。该项锚定的是前史条件与信息背景，不把条件直接写成唯一因果。来源保留：官修编年和地方材料的记录密度与立场并不相同；本行不将条件、传闻、呈报或后见解释直接等同为确定因果。
\eventanchor{EQ-1729-JUNJICHU-FOUNDING-DECISION}{围绕「雍正帝设军机房，后来发…}雍正七年（1729年），围绕「雍正帝设军机房，后来发展为军机处，处理紧急军政文书」的命令、调度、照会或呈报形成独立的中央—地方行动文书链。核心取证：《清世宗实录》雍正七年相关条；宫中朱批奏折、内阁大库题本与《清史稿》相关志传。该项锚定的是命令或决策环节，命令抵达和结果仍须另外核验。来源保留：奏折、上谕、部议、军机档或条约可证明行政动作和机构立场，不自动证明所有地区、群体或对手按其意图行动。
\eventanchor{EQ-1729-JUNJICHU-FOUNDING-AFTERMATH}{「雍正帝设军机房，后来发展为…}雍正七年（1729年），「雍正帝设军机房，后来发展为军机处，处理紧急军政文书」引发的善后、执行或地方回应构成可由不同材料追踪的后续节点。核心取证：《清世宗实录》雍正七年相关条；宫中朱批奏折、内阁大库题本与《清史稿》相关志传；相关地区的档案、志书、契约、报刊或对方材料。该项锚定的是后续追踪问题，影响范围须按地区与群体分别判断。来源保留：战后、改革后或条约后的记忆、地方记录和官方总结常具有事后选择性；后续影响须按地点、群体和时间分层，不作整体化推断。
\subsection{1730—1739}
\eventanchor{EQ-1730-SOUTHWEST-MIAO-GOVERNANCE-PRELUDE}{「清廷围绕苗疆设置、屯兵与地…}雍正八年（1730年），「清廷围绕苗疆设置、屯兵与地方纠纷加强治理」之前的条件、信息和争议已通过中央与地方材料显现，构成该事件可追踪的前史节点。核心取证：《清世宗实录》雍正八年相关条；军机处档、理藩院档或相关方略；同时期地方档案、地方志、日记或对方材料应按具体地区补检。该项锚定的是前史条件与信息背景，不把条件直接写成唯一因果。来源保留：官修编年和地方材料的记录密度与立场并不相同；本行不将条件、传闻、呈报或后见解释直接等同为确定因果。
\eventanchor{EQ-1730-SOUTHWEST-MIAO-GOVERNANCE-DECISION}{围绕「清廷围绕苗疆设置、屯兵…}雍正八年（1730年），围绕「清廷围绕苗疆设置、屯兵与地方纠纷加强治理」的命令、调度、照会或呈报形成独立的中央—地方行动文书链。核心取证：《清世宗实录》雍正八年相关条；军机处档、理藩院档或相关方略。该项锚定的是命令或决策环节，命令抵达和结果仍须另外核验。来源保留：奏折、上谕、部议、军机档或条约可证明行政动作和机构立场，不自动证明所有地区、群体或对手按其意图行动。
\eventanchor{EQ-1730-SOUTHWEST-MIAO-GOVERNANCE-AFTERMATH}{「清廷围绕苗疆设置、屯兵与地…}雍正八年（1730年），「清廷围绕苗疆设置、屯兵与地方纠纷加强治理」引发的善后、执行或地方回应构成可由不同材料追踪的后续节点。核心取证：《清世宗实录》雍正八年相关条；军机处档、理藩院档或相关方略；相关地区的档案、志书、契约、报刊或对方材料。该项锚定的是后续追踪问题，影响范围须按地区与群体分别判断。来源保留：战后、改革后或条约后的记忆、地方记录和官方总结常具有事后选择性；后续影响须按地点、群体和时间分层，不作整体化推断。
\eventanchor{EQ-1731-Dzungar-WAR-REVERSE-PRELUDE}{「清军在对准噶尔战争中受挫，…}雍正九年（1731年），「清军在对准噶尔战争中受挫，西北远征暴露补给与指挥困难」之前的条件、信息和争议已通过中央与地方材料显现，构成该事件可追踪的前史节点。核心取证：《清世宗实录》雍正九年相关条；军机处档、战事方略及地方奏报；同时期地方档案、地方志、日记或对方材料应按具体地区补检。该项锚定的是前史条件与信息背景，不把条件直接写成唯一因果。来源保留：官修编年和地方材料的记录密度与立场并不相同；本行不将条件、传闻、呈报或后见解释直接等同为确定因果。
\eventanchor{EQ-1731-Dzungar-WAR-REVERSE-DECISION}{围绕「清军在对准噶尔战争中受…}雍正九年（1731年），围绕「清军在对准噶尔战争中受挫，西北远征暴露补给与指挥困难」的命令、调度、照会或呈报形成独立的中央—地方行动文书链。核心取证：《清世宗实录》雍正九年相关条；军机处档、战事方略及地方奏报。该项锚定的是命令或决策环节，命令抵达和结果仍须另外核验。来源保留：奏折、上谕、部议、军机档或条约可证明行政动作和机构立场，不自动证明所有地区、群体或对手按其意图行动。
\eventanchor{EQ-1731-Dzungar-WAR-REVERSE-AFTERMATH}{「清军在对准噶尔战争中受挫，…}雍正九年（1731年），「清军在对准噶尔战争中受挫，西北远征暴露补给与指挥困难」引发的善后、执行或地方回应构成可由不同材料追踪的后续节点。核心取证：《清世宗实录》雍正九年相关条；军机处档、战事方略及地方奏报；相关地区的档案、志书、契约、报刊或对方材料。该项锚定的是后续追踪问题，影响范围须按地区与群体分别判断。来源保留：战后、改革后或条约后的记忆、地方记录和官方总结常具有事后选择性；后续影响须按地点、群体和时间分层，不作整体化推断。
\eventanchor{EQ-1732-KOBDO-BATTLE-PRELUDE}{「清军在和通泊等战事失利，西…}雍正十年（1732年），「清军在和通泊等战事失利，西北战局一度紧张」之前的条件、信息和争议已通过中央与地方材料显现，构成该事件可追踪的前史节点。核心取证：《清世宗实录》雍正十年相关条；军机处档、战事方略及地方奏报；同时期地方档案、地方志、日记或对方材料应按具体地区补检。该项锚定的是前史条件与信息背景，不把条件直接写成唯一因果。来源保留：官修编年和地方材料的记录密度与立场并不相同；本行不将条件、传闻、呈报或后见解释直接等同为确定因果。
\eventanchor{EQ-1732-KOBDO-BATTLE-DECISION}{围绕「清军在和通泊等战事失利…}雍正十年（1732年），围绕「清军在和通泊等战事失利，西北战局一度紧张」的命令、调度、照会或呈报形成独立的中央—地方行动文书链。核心取证：《清世宗实录》雍正十年相关条；军机处档、战事方略及地方奏报。该项锚定的是命令或决策环节，命令抵达和结果仍须另外核验。来源保留：奏折、上谕、部议、军机档或条约可证明行政动作和机构立场，不自动证明所有地区、群体或对手按其意图行动。
\eventanchor{EQ-1732-KOBDO-BATTLE-AFTERMATH}{「清军在和通泊等战事失利，西…}雍正十年（1732年），「清军在和通泊等战事失利，西北战局一度紧张」引发的善后、执行或地方回应构成可由不同材料追踪的后续节点。核心取证：《清世宗实录》雍正十年相关条；军机处档、战事方略及地方奏报；相关地区的档案、志书、契约、报刊或对方材料。该项锚定的是后续追踪问题，影响范围须按地区与群体分别判断。来源保留：战后、改革后或条约后的记忆、地方记录和官方总结常具有事后选择性；后续影响须按地点、群体和时间分层，不作整体化推断。
\eventanchor{EQ-1733-YONGZHENG-FISCAL-INSPECTION-PRELUDE}{「雍正政府持续以奏销、亏空追…}雍正十一年（1733年），「雍正政府持续以奏销、亏空追查和官员考成为手段整顿地方财政」之前的条件、信息和争议已通过中央与地方材料显现，构成该事件可追踪的前史节点。核心取证：《清世宗实录》雍正十一年相关条；户部奏销、盐政、河工或海关档案；同时期地方档案、地方志、日记或对方材料应按具体地区补检。该项锚定的是前史条件与信息背景，不把条件直接写成唯一因果。来源保留：官修编年和地方材料的记录密度与立场并不相同；本行不将条件、传闻、呈报或后见解释直接等同为确定因果。
\eventanchor{EQ-1733-YONGZHENG-FISCAL-INSPECTION-DECISION}{围绕「雍正政府持续以奏销、亏…}雍正十一年（1733年），围绕「雍正政府持续以奏销、亏空追查和官员考成为手段整顿地方财政」的命令、调度、照会或呈报形成独立的中央—地方行动文书链。核心取证：《清世宗实录》雍正十一年相关条；户部奏销、盐政、河工或海关档案。该项锚定的是命令或决策环节，命令抵达和结果仍须另外核验。来源保留：奏折、上谕、部议、军机档或条约可证明行政动作和机构立场，不自动证明所有地区、群体或对手按其意图行动。
\eventanchor{EQ-1733-YONGZHENG-FISCAL-INSPECTION-AFTERMATH}{「雍正政府持续以奏销、亏空追…}雍正十一年（1733年），「雍正政府持续以奏销、亏空追查和官员考成为手段整顿地方财政」引发的善后、执行或地方回应构成可由不同材料追踪的后续节点。核心取证：《清世宗实录》雍正十一年相关条；户部奏销、盐政、河工或海关档案；相关地区的档案、志书、契约、报刊或对方材料。该项锚定的是后续追踪问题，影响范围须按地区与群体分别判断。来源保留：战后、改革后或条约后的记忆、地方记录和官方总结常具有事后选择性；后续影响须按地点、群体和时间分层，不作整体化推断。
\eventanchor{EQ-1734-CHRISTIAN-ENFORCEMENT-PRELUDE}{「地方官继续执行限制天主教的…}雍正十二年（1734年），「地方官继续执行限制天主教的政策，传教士和教民的活动受地域差异影响」之前的条件、信息和争议已通过中央与地方材料显现，构成该事件可追踪的前史节点。核心取证：《清世宗实录》雍正十二年相关条；地方志、督抚奏折及地方档案；同时期地方档案、地方志、日记或对方材料应按具体地区补检。该项锚定的是前史条件与信息背景，不把条件直接写成唯一因果。来源保留：官修编年和地方材料的记录密度与立场并不相同；本行不将条件、传闻、呈报或后见解释直接等同为确定因果。
\eventanchor{EQ-1734-CHRISTIAN-ENFORCEMENT-DECISION}{围绕「地方官继续执行限制天主…}雍正十二年（1734年），围绕「地方官继续执行限制天主教的政策，传教士和教民的活动受地域差异影响」的命令、调度、照会或呈报形成独立的中央—地方行动文书链。核心取证：《清世宗实录》雍正十二年相关条；地方志、督抚奏折及地方档案。该项锚定的是命令或决策环节，命令抵达和结果仍须另外核验。来源保留：奏折、上谕、部议、军机档或条约可证明行政动作和机构立场，不自动证明所有地区、群体或对手按其意图行动。
\eventanchor{EQ-1734-CHRISTIAN-ENFORCEMENT-AFTERMATH}{「地方官继续执行限制天主教的…}雍正十二年（1734年），「地方官继续执行限制天主教的政策，传教士和教民的活动受地域差异影响」引发的善后、执行或地方回应构成可由不同材料追踪的后续节点。核心取证：《清世宗实录》雍正十二年相关条；地方志、督抚奏折及地方档案；相关地区的档案、志书、契约、报刊或对方材料。该项锚定的是后续追踪问题，影响范围须按地区与群体分别判断。来源保留：战后、改革后或条约后的记忆、地方记录和官方总结常具有事后选择性；后续影响须按地点、群体和时间分层，不作整体化推断。
\eventanchor{EQ-1735-QIANLONG-ACCESSION-PRELUDE}{「乾隆帝即位，雍正时期形成的…}雍正十三年（1735年），「乾隆帝即位，雍正时期形成的财政与军机决策机制进入新朝」之前的条件、信息和争议已通过中央与地方材料显现，构成该事件可追踪的前史节点。核心取证：《清世宗实录》雍正十三年相关条；宫中朱批奏折、内阁大库题本与《清史稿》相关志传；同时期地方档案、地方志、日记或对方材料应按具体地区补检。该项锚定的是前史条件与信息背景，不把条件直接写成唯一因果。来源保留：官修编年和地方材料的记录密度与立场并不相同；本行不将条件、传闻、呈报或后见解释直接等同为确定因果。
\eventanchor{EQ-1735-QIANLONG-ACCESSION-DECISION}{围绕「乾隆帝即位，雍正时期形…}雍正十三年（1735年），围绕「乾隆帝即位，雍正时期形成的财政与军机决策机制进入新朝」的命令、调度、照会或呈报形成独立的中央—地方行动文书链。核心取证：《清世宗实录》雍正十三年相关条；宫中朱批奏折、内阁大库题本与《清史稿》相关志传。该项锚定的是命令或决策环节，命令抵达和结果仍须另外核验。来源保留：奏折、上谕、部议、军机档或条约可证明行政动作和机构立场，不自动证明所有地区、群体或对手按其意图行动。
\eventanchor{EQ-1735-QIANLONG-ACCESSION-AFTERMATH}{「乾隆帝即位，雍正时期形成的…}雍正十三年（1735年），「乾隆帝即位，雍正时期形成的财政与军机决策机制进入新朝」引发的善后、执行或地方回应构成可由不同材料追踪的后续节点。核心取证：《清世宗实录》雍正十三年相关条；宫中朱批奏折、内阁大库题本与《清史稿》相关志传；相关地区的档案、志书、契约、报刊或对方材料。该项锚定的是后续追踪问题，影响范围须按地区与群体分别判断。来源保留：战后、改革后或条约后的记忆、地方记录和官方总结常具有事后选择性；后续影响须按地点、群体和时间分层，不作整体化推断。
\eventanchor{EQ-1736-AMNESTY-AND-RECTIFICATION-PRELUDE}{「乾隆初年颁布恩诏并调整雍正…}乾隆元年（1736年），「乾隆初年颁布恩诏并调整雍正后期政治案件的处置」之前的条件、信息和争议已通过中央与地方材料显现，构成该事件可追踪的前史节点。核心取证：《清高宗实录》乾隆元年相关条；宫中朱批奏折、内阁大库题本与《清史稿》相关志传；同时期地方档案、地方志、日记或对方材料应按具体地区补检。该项锚定的是前史条件与信息背景，不把条件直接写成唯一因果。来源保留：官修编年和地方材料的记录密度与立场并不相同；本行不将条件、传闻、呈报或后见解释直接等同为确定因果。
\eventanchor{EQ-1736-AMNESTY-AND-RECTIFICATION-DECISION}{围绕「乾隆初年颁布恩诏并调整…}乾隆元年（1736年），围绕「乾隆初年颁布恩诏并调整雍正后期政治案件的处置」的命令、调度、照会或呈报形成独立的中央—地方行动文书链。核心取证：《清高宗实录》乾隆元年相关条；宫中朱批奏折、内阁大库题本与《清史稿》相关志传。该项锚定的是命令或决策环节，命令抵达和结果仍须另外核验。来源保留：奏折、上谕、部议、军机档或条约可证明行政动作和机构立场，不自动证明所有地区、群体或对手按其意图行动。
\eventanchor{EQ-1736-AMNESTY-AND-RECTIFICATION-AFTERMATH}{「乾隆初年颁布恩诏并调整雍正…}乾隆元年（1736年），「乾隆初年颁布恩诏并调整雍正后期政治案件的处置」引发的善后、执行或地方回应构成可由不同材料追踪的后续节点。核心取证：《清高宗实录》乾隆元年相关条；宫中朱批奏折、内阁大库题本与《清史稿》相关志传；相关地区的档案、志书、契约、报刊或对方材料。该项锚定的是后续追踪问题，影响范围须按地区与群体分别判断。来源保留：战后、改革后或条约后的记忆、地方记录和官方总结常具有事后选择性；后续影响须按地点、群体和时间分层，不作整体化推断。
\eventanchor{EQ-1737-Dzungar-NEGOTIATION-PRELUDE}{「清廷在战争压力下继续处理与…}乾隆二年（1737年），「清廷在战争压力下继续处理与准噶尔的使节、盟旗和军事情报」之前的条件、信息和争议已通过中央与地方材料显现，构成该事件可追踪的前史节点。核心取证：《清高宗实录》乾隆二年相关条；《筹办夷务始末》相关年卷与中外照会、条约文本；同时期地方档案、地方志、日记或对方材料应按具体地区补检。该项锚定的是前史条件与信息背景，不把条件直接写成唯一因果。来源保留：官修编年和地方材料的记录密度与立场并不相同；本行不将条件、传闻、呈报或后见解释直接等同为确定因果。
\eventanchor{EQ-1737-Dzungar-NEGOTIATION-DECISION}{围绕「清廷在战争压力下继续处…}乾隆二年（1737年），围绕「清廷在战争压力下继续处理与准噶尔的使节、盟旗和军事情报」的命令、调度、照会或呈报形成独立的中央—地方行动文书链。核心取证：《清高宗实录》乾隆二年相关条；《筹办夷务始末》相关年卷与中外照会、条约文本。该项锚定的是命令或决策环节，命令抵达和结果仍须另外核验。来源保留：奏折、上谕、部议、军机档或条约可证明行政动作和机构立场，不自动证明所有地区、群体或对手按其意图行动。
\eventanchor{EQ-1737-Dzungar-NEGOTIATION-AFTERMATH}{「清廷在战争压力下继续处理与…}乾隆二年（1737年），「清廷在战争压力下继续处理与准噶尔的使节、盟旗和军事情报」引发的善后、执行或地方回应构成可由不同材料追踪的后续节点。核心取证：《清高宗实录》乾隆二年相关条；《筹办夷务始末》相关年卷与中外照会、条约文本；相关地区的档案、志书、契约、报刊或对方材料。该项锚定的是后续追踪问题，影响范围须按地区与群体分别判断。来源保留：战后、改革后或条约后的记忆、地方记录和官方总结常具有事后选择性；后续影响须按地点、群体和时间分层，不作整体化推断。
\eventanchor{EQ-1738-FIRST-JINCHUAN-WAR-PRELUDE}{「清军开始大规模干预大金川，…}乾隆三年（1738年），「清军开始大规模干预大金川，山地战与土司关系限制军事推进」之前的条件、信息和争议已通过中央与地方材料显现，构成该事件可追踪的前史节点。核心取证：《清高宗实录》乾隆三年相关条；军机处档、战事方略及地方奏报；同时期地方档案、地方志、日记或对方材料应按具体地区补检。该项锚定的是前史条件与信息背景，不把条件直接写成唯一因果。来源保留：官修编年和地方材料的记录密度与立场并不相同；本行不将条件、传闻、呈报或后见解释直接等同为确定因果。
\eventanchor{EQ-1738-FIRST-JINCHUAN-WAR-DECISION}{围绕「清军开始大规模干预大金…}乾隆三年（1738年），围绕「清军开始大规模干预大金川，山地战与土司关系限制军事推进」的命令、调度、照会或呈报形成独立的中央—地方行动文书链。核心取证：《清高宗实录》乾隆三年相关条；军机处档、战事方略及地方奏报。该项锚定的是命令或决策环节，命令抵达和结果仍须另外核验。来源保留：奏折、上谕、部议、军机档或条约可证明行政动作和机构立场，不自动证明所有地区、群体或对手按其意图行动。
\eventanchor{EQ-1738-FIRST-JINCHUAN-WAR-AFTERMATH}{「清军开始大规模干预大金川，…}乾隆三年（1738年），「清军开始大规模干预大金川，山地战与土司关系限制军事推进」引发的善后、执行或地方回应构成可由不同材料追踪的后续节点。核心取证：《清高宗实录》乾隆三年相关条；军机处档、战事方略及地方奏报；相关地区的档案、志书、契约、报刊或对方材料。该项锚定的是后续追踪问题，影响范围须按地区与群体分别判断。来源保留：战后、改革后或条约后的记忆、地方记录和官方总结常具有事后选择性；后续影响须按地点、群体和时间分层，不作整体化推断。
\subsection{1740—1749}
\eventanchor{EQ-1740-NORTH-CHINA-FAMINE-PRELUDE}{「华北旱灾和粮价问题促使地方…}乾隆五年（1740年），「华北旱灾和粮价问题促使地方奏报与赈济措施增加」之前的条件、信息和争议已通过中央与地方材料显现，构成该事件可追踪的前史节点。核心取证：《清高宗实录》乾隆五年相关条；地方志、督抚奏折及地方档案；同时期地方档案、地方志、日记或对方材料应按具体地区补检。该项锚定的是前史条件与信息背景，不把条件直接写成唯一因果。来源保留：官修编年和地方材料的记录密度与立场并不相同；本行不将条件、传闻、呈报或后见解释直接等同为确定因果。
\eventanchor{EQ-1740-NORTH-CHINA-FAMINE-DECISION}{围绕「华北旱灾和粮价问题促使…}乾隆五年（1740年），围绕「华北旱灾和粮价问题促使地方奏报与赈济措施增加」的命令、调度、照会或呈报形成独立的中央—地方行动文书链。核心取证：《清高宗实录》乾隆五年相关条；地方志、督抚奏折及地方档案。该项锚定的是命令或决策环节，命令抵达和结果仍须另外核验。来源保留：奏折、上谕、部议、军机档或条约可证明行政动作和机构立场，不自动证明所有地区、群体或对手按其意图行动。
\eventanchor{EQ-1740-NORTH-CHINA-FAMINE-AFTERMATH}{「华北旱灾和粮价问题促使地方…}乾隆五年（1740年），「华北旱灾和粮价问题促使地方奏报与赈济措施增加」引发的善后、执行或地方回应构成可由不同材料追踪的后续节点。核心取证：《清高宗实录》乾隆五年相关条；地方志、督抚奏折及地方档案；相关地区的档案、志书、契约、报刊或对方材料。该项锚定的是后续追踪问题，影响范围须按地区与群体分别判断。来源保留：战后、改革后或条约后的记忆、地方记录和官方总结常具有事后选择性；后续影响须按地点、群体和时间分层，不作整体化推断。
\eventanchor{EQ-1741-MIAO-FRONTIER-CAMPAIGN-PRELUDE}{「清廷在西南继续以驻兵、设治…}乾隆六年（1741年），「清廷在西南继续以驻兵、设治和军事行动处理苗疆冲突」之前的条件、信息和争议已通过中央与地方材料显现，构成该事件可追踪的前史节点。核心取证：《清高宗实录》乾隆六年相关条；军机处档、理藩院档或相关方略；同时期地方档案、地方志、日记或对方材料应按具体地区补检。该项锚定的是前史条件与信息背景，不把条件直接写成唯一因果。来源保留：官修编年和地方材料的记录密度与立场并不相同；本行不将条件、传闻、呈报或后见解释直接等同为确定因果。
\eventanchor{EQ-1741-MIAO-FRONTIER-CAMPAIGN-DECISION}{围绕「清廷在西南继续以驻兵、…}乾隆六年（1741年），围绕「清廷在西南继续以驻兵、设治和军事行动处理苗疆冲突」的命令、调度、照会或呈报形成独立的中央—地方行动文书链。核心取证：《清高宗实录》乾隆六年相关条；军机处档、理藩院档或相关方略。该项锚定的是命令或决策环节，命令抵达和结果仍须另外核验。来源保留：奏折、上谕、部议、军机档或条约可证明行政动作和机构立场，不自动证明所有地区、群体或对手按其意图行动。
\eventanchor{EQ-1741-MIAO-FRONTIER-CAMPAIGN-AFTERMATH}{「清廷在西南继续以驻兵、设治…}乾隆六年（1741年），「清廷在西南继续以驻兵、设治和军事行动处理苗疆冲突」引发的善后、执行或地方回应构成可由不同材料追踪的后续节点。核心取证：《清高宗实录》乾隆六年相关条；军机处档、理藩院档或相关方略；相关地区的档案、志书、契约、报刊或对方材料。该项锚定的是后续追踪问题，影响范围须按地区与群体分别判断。来源保留：战后、改革后或条约后的记忆、地方记录和官方总结常具有事后选择性；后续影响须按地点、群体和时间分层，不作整体化推断。
\eventanchor{EQ-1742-BANNER-EXIT-ORDER-PRELUDE}{「清廷允许部分汉军旗人出旗，…}乾隆七年（1742年），「清廷允许部分汉军旗人出旗，调整旗籍与生计安排」之前的条件、信息和争议已通过中央与地方材料显现，构成该事件可追踪的前史节点。核心取证：《清高宗实录》乾隆七年相关条；地方志、督抚奏折及地方档案；同时期地方档案、地方志、日记或对方材料应按具体地区补检。该项锚定的是前史条件与信息背景，不把条件直接写成唯一因果。来源保留：官修编年和地方材料的记录密度与立场并不相同；本行不将条件、传闻、呈报或后见解释直接等同为确定因果。
\eventanchor{EQ-1742-BANNER-EXIT-ORDER-DECISION}{围绕「清廷允许部分汉军旗人出…}乾隆七年（1742年），围绕「清廷允许部分汉军旗人出旗，调整旗籍与生计安排」的命令、调度、照会或呈报形成独立的中央—地方行动文书链。核心取证：《清高宗实录》乾隆七年相关条；地方志、督抚奏折及地方档案。该项锚定的是命令或决策环节，命令抵达和结果仍须另外核验。来源保留：奏折、上谕、部议、军机档或条约可证明行政动作和机构立场，不自动证明所有地区、群体或对手按其意图行动。
\eventanchor{EQ-1742-BANNER-EXIT-ORDER-AFTERMATH}{「清廷允许部分汉军旗人出旗，…}乾隆七年（1742年），「清廷允许部分汉军旗人出旗，调整旗籍与生计安排」引发的善后、执行或地方回应构成可由不同材料追踪的后续节点。核心取证：《清高宗实录》乾隆七年相关条；地方志、督抚奏折及地方档案；相关地区的档案、志书、契约、报刊或对方材料。该项锚定的是后续追踪问题，影响范围须按地区与群体分别判断。来源保留：战后、改革后或条约后的记忆、地方记录和官方总结常具有事后选择性；后续影响须按地点、群体和时间分层，不作整体化推断。
\eventanchor{EQ-1744-JESUIT-CALENDAR-SERVICE-PRELUDE}{「宫廷继续利用部分西洋传教士…}乾隆九年（1744年），「宫廷继续利用部分西洋传教士的历算与技术服务，同时维持宗教限制」之前的条件、信息和争议已通过中央与地方材料显现，构成该事件可追踪的前史节点。核心取证：《清高宗实录》乾隆九年相关条；内阁档、文集或报刊相关记录；同时期地方档案、地方志、日记或对方材料应按具体地区补检。该项锚定的是前史条件与信息背景，不把条件直接写成唯一因果。来源保留：官修编年和地方材料的记录密度与立场并不相同；本行不将条件、传闻、呈报或后见解释直接等同为确定因果。
\eventanchor{EQ-1744-JESUIT-CALENDAR-SERVICE-DECISION}{围绕「宫廷继续利用部分西洋传…}乾隆九年（1744年），围绕「宫廷继续利用部分西洋传教士的历算与技术服务，同时维持宗教限制」的命令、调度、照会或呈报形成独立的中央—地方行动文书链。核心取证：《清高宗实录》乾隆九年相关条；内阁档、文集或报刊相关记录。该项锚定的是命令或决策环节，命令抵达和结果仍须另外核验。来源保留：奏折、上谕、部议、军机档或条约可证明行政动作和机构立场，不自动证明所有地区、群体或对手按其意图行动。
\eventanchor{EQ-1744-JESUIT-CALENDAR-SERVICE-AFTERMATH}{「宫廷继续利用部分西洋传教士…}乾隆九年（1744年），「宫廷继续利用部分西洋传教士的历算与技术服务，同时维持宗教限制」引发的善后、执行或地方回应构成可由不同材料追踪的后续节点。核心取证：《清高宗实录》乾隆九年相关条；内阁档、文集或报刊相关记录；相关地区的档案、志书、契约、报刊或对方材料。该项锚定的是后续追踪问题，影响范围须按地区与群体分别判断。来源保留：战后、改革后或条约后的记忆、地方记录和官方总结常具有事后选择性；后续影响须按地点、群体和时间分层，不作整体化推断。
\eventanchor{EQ-1745-JINCHUAN-NEGOTIATION-PRELUDE}{「清廷在大金川战局中尝试通过…}乾隆十年（1745年），「清廷在大金川战局中尝试通过招抚、盟约和军事压力重组地方关系」之前的条件、信息和争议已通过中央与地方材料显现，构成该事件可追踪的前史节点。核心取证：《清高宗实录》乾隆十年相关条；军机处档、理藩院档或相关方略；同时期地方档案、地方志、日记或对方材料应按具体地区补检。该项锚定的是前史条件与信息背景，不把条件直接写成唯一因果。来源保留：官修编年和地方材料的记录密度与立场并不相同；本行不将条件、传闻、呈报或后见解释直接等同为确定因果。
\eventanchor{EQ-1745-JINCHUAN-NEGOTIATION-DECISION}{围绕「清廷在大金川战局中尝试…}乾隆十年（1745年），围绕「清廷在大金川战局中尝试通过招抚、盟约和军事压力重组地方关系」的命令、调度、照会或呈报形成独立的中央—地方行动文书链。核心取证：《清高宗实录》乾隆十年相关条；军机处档、理藩院档或相关方略。该项锚定的是命令或决策环节，命令抵达和结果仍须另外核验。来源保留：奏折、上谕、部议、军机档或条约可证明行政动作和机构立场，不自动证明所有地区、群体或对手按其意图行动。
\eventanchor{EQ-1745-JINCHUAN-NEGOTIATION-AFTERMATH}{「清廷在大金川战局中尝试通过…}乾隆十年（1745年），「清廷在大金川战局中尝试通过招抚、盟约和军事压力重组地方关系」引发的善后、执行或地方回应构成可由不同材料追踪的后续节点。核心取证：《清高宗实录》乾隆十年相关条；军机处档、理藩院档或相关方略；相关地区的档案、志书、契约、报刊或对方材料。该项锚定的是后续追踪问题，影响范围须按地区与群体分别判断。来源保留：战后、改革后或条约后的记忆、地方记录和官方总结常具有事后选择性；后续影响须按地点、群体和时间分层，不作整体化推断。
\eventanchor{EQ-1746-FIRST-JINCHUAN-SETTLEMENT-PRELUDE}{「第一次金川战争以和议和行政…}乾隆十一年（1746年），「第一次金川战争以和议和行政安排告一段落」之前的条件、信息和争议已通过中央与地方材料显现，构成该事件可追踪的前史节点。核心取证：《清高宗实录》乾隆十一年相关条；军机处档、战事方略及地方奏报；同时期地方档案、地方志、日记或对方材料应按具体地区补检。该项锚定的是前史条件与信息背景，不把条件直接写成唯一因果。来源保留：官修编年和地方材料的记录密度与立场并不相同；本行不将条件、传闻、呈报或后见解释直接等同为确定因果。
\eventanchor{EQ-1746-FIRST-JINCHUAN-SETTLEMENT-DECISION}{围绕「第一次金川战争以和议和…}乾隆十一年（1746年），围绕「第一次金川战争以和议和行政安排告一段落」的命令、调度、照会或呈报形成独立的中央—地方行动文书链。核心取证：《清高宗实录》乾隆十一年相关条；军机处档、战事方略及地方奏报。该项锚定的是命令或决策环节，命令抵达和结果仍须另外核验。来源保留：奏折、上谕、部议、军机档或条约可证明行政动作和机构立场，不自动证明所有地区、群体或对手按其意图行动。
\eventanchor{EQ-1746-FIRST-JINCHUAN-SETTLEMENT-AFTERMATH}{「第一次金川战争以和议和行政…}乾隆十一年（1746年），「第一次金川战争以和议和行政安排告一段落」引发的善后、执行或地方回应构成可由不同材料追踪的后续节点。核心取证：《清高宗实录》乾隆十一年相关条；军机处档、战事方略及地方奏报；相关地区的档案、志书、契约、报刊或对方材料。该项锚定的是后续追踪问题，影响范围须按地区与群体分别判断。来源保留：战后、改革后或条约后的记忆、地方记录和官方总结常具有事后选择性；后续影响须按地点、群体和时间分层，不作整体化推断。
\eventanchor{EQ-1748-SONG-AND-CASE-PRELUDE}{「乾隆朝通过大案处分官员，强…}乾隆十三年（1748年），「乾隆朝通过大案处分官员，强化对地方文书、钱粮和吏治的监督」之前的条件、信息和争议已通过中央与地方材料显现，构成该事件可追踪的前史节点。核心取证：《清高宗实录》乾隆十三年相关条；宫中朱批奏折、内阁大库题本与《清史稿》相关志传；同时期地方档案、地方志、日记或对方材料应按具体地区补检。该项锚定的是前史条件与信息背景，不把条件直接写成唯一因果。来源保留：官修编年和地方材料的记录密度与立场并不相同；本行不将条件、传闻、呈报或后见解释直接等同为确定因果。
\eventanchor{EQ-1748-SONG-AND-CASE-DECISION}{围绕「乾隆朝通过大案处分官员…}乾隆十三年（1748年），围绕「乾隆朝通过大案处分官员，强化对地方文书、钱粮和吏治的监督」的命令、调度、照会或呈报形成独立的中央—地方行动文书链。核心取证：《清高宗实录》乾隆十三年相关条；宫中朱批奏折、内阁大库题本与《清史稿》相关志传。该项锚定的是命令或决策环节，命令抵达和结果仍须另外核验。来源保留：奏折、上谕、部议、军机档或条约可证明行政动作和机构立场，不自动证明所有地区、群体或对手按其意图行动。
\eventanchor{EQ-1748-SONG-AND-CASE-AFTERMATH}{「乾隆朝通过大案处分官员，强…}乾隆十三年（1748年），「乾隆朝通过大案处分官员，强化对地方文书、钱粮和吏治的监督」引发的善后、执行或地方回应构成可由不同材料追踪的后续节点。核心取证：《清高宗实录》乾隆十三年相关条；宫中朱批奏折、内阁大库题本与《清史稿》相关志传；相关地区的档案、志书、契约、报刊或对方材料。该项锚定的是后续追踪问题，影响范围须按地区与群体分别判断。来源保留：战后、改革后或条约后的记忆、地方记录和官方总结常具有事后选择性；后续影响须按地点、群体和时间分层，不作整体化推断。
\eventanchor{EQ-1749-JIANGNAN-TAX-REVIEW-PRELUDE}{「江南钱粮、漕运和士绅关系继…}乾隆十四年（1749年），「江南钱粮、漕运和士绅关系继续在户部与地方奏折中被讨论」之前的条件、信息和争议已通过中央与地方材料显现，构成该事件可追踪的前史节点。核心取证：《清高宗实录》乾隆十四年相关条；户部奏销、盐政、河工或海关档案；同时期地方档案、地方志、日记或对方材料应按具体地区补检。该项锚定的是前史条件与信息背景，不把条件直接写成唯一因果。来源保留：官修编年和地方材料的记录密度与立场并不相同；本行不将条件、传闻、呈报或后见解释直接等同为确定因果。
\eventanchor{EQ-1749-JIANGNAN-TAX-REVIEW-DECISION}{围绕「江南钱粮、漕运和士绅关…}乾隆十四年（1749年），围绕「江南钱粮、漕运和士绅关系继续在户部与地方奏折中被讨论」的命令、调度、照会或呈报形成独立的中央—地方行动文书链。核心取证：《清高宗实录》乾隆十四年相关条；户部奏销、盐政、河工或海关档案。该项锚定的是命令或决策环节，命令抵达和结果仍须另外核验。来源保留：奏折、上谕、部议、军机档或条约可证明行政动作和机构立场，不自动证明所有地区、群体或对手按其意图行动。
\eventanchor{EQ-1749-JIANGNAN-TAX-REVIEW-AFTERMATH}{「江南钱粮、漕运和士绅关系继…}乾隆十四年（1749年），「江南钱粮、漕运和士绅关系继续在户部与地方奏折中被讨论」引发的善后、执行或地方回应构成可由不同材料追踪的后续节点。核心取证：《清高宗实录》乾隆十四年相关条；户部奏销、盐政、河工或海关档案；相关地区的档案、志书、契约、报刊或对方材料。该项锚定的是后续追踪问题，影响范围须按地区与群体分别判断。来源保留：战后、改革后或条约后的记忆、地方记录和官方总结常具有事后选择性；后续影响须按地点、群体和时间分层，不作整体化推断。
\subsection{1750—1759}
\eventanchor{EQ-1750-LHASA-UPRISING-PRELUDE}{「拉萨发生政治冲突，驻藏大臣…}乾隆十五年（1750年），「拉萨发生政治冲突，驻藏大臣和清军介入处置」之前的条件、信息和争议已通过中央与地方材料显现，构成该事件可追踪的前史节点。核心取证：《清高宗实录》乾隆十五年相关条；军机处档、理藩院档或相关方略；同时期地方档案、地方志、日记或对方材料应按具体地区补检。该项锚定的是前史条件与信息背景，不把条件直接写成唯一因果。来源保留：官修编年和地方材料的记录密度与立场并不相同；本行不将条件、传闻、呈报或后见解释直接等同为确定因果。
\eventanchor{EQ-1750-LHASA-UPRISING-DECISION}{围绕「拉萨发生政治冲突，驻藏…}乾隆十五年（1750年），围绕「拉萨发生政治冲突，驻藏大臣和清军介入处置」的命令、调度、照会或呈报形成独立的中央—地方行动文书链。核心取证：《清高宗实录》乾隆十五年相关条；军机处档、理藩院档或相关方略。该项锚定的是命令或决策环节，命令抵达和结果仍须另外核验。来源保留：奏折、上谕、部议、军机档或条约可证明行政动作和机构立场，不自动证明所有地区、群体或对手按其意图行动。
\eventanchor{EQ-1750-LHASA-UPRISING-AFTERMATH}{「拉萨发生政治冲突，驻藏大臣…}乾隆十五年（1750年），「拉萨发生政治冲突，驻藏大臣和清军介入处置」引发的善后、执行或地方回应构成可由不同材料追踪的后续节点。核心取证：《清高宗实录》乾隆十五年相关条；军机处档、理藩院档或相关方略；相关地区的档案、志书、契约、报刊或对方材料。该项锚定的是后续追踪问题，影响范围须按地区与群体分别判断。来源保留：战后、改革后或条约后的记忆、地方记录和官方总结常具有事后选择性；后续影响须按地点、群体和时间分层，不作整体化推断。
\eventanchor{EQ-1751-TIBET-ORDINANCE-PRELUDE}{「清廷在拉萨事变后重整驻藏大…}乾隆十六年（1751年），「清廷在拉萨事变后重整驻藏大臣和噶厦相关制度安排」之前的条件、信息和争议已通过中央与地方材料显现，构成该事件可追踪的前史节点。核心取证：《清高宗实录》乾隆十六年相关条；军机处档、理藩院档或相关方略；同时期地方档案、地方志、日记或对方材料应按具体地区补检。该项锚定的是前史条件与信息背景，不把条件直接写成唯一因果。来源保留：官修编年和地方材料的记录密度与立场并不相同；本行不将条件、传闻、呈报或后见解释直接等同为确定因果。
\eventanchor{EQ-1751-TIBET-ORDINANCE-DECISION}{围绕「清廷在拉萨事变后重整驻…}乾隆十六年（1751年），围绕「清廷在拉萨事变后重整驻藏大臣和噶厦相关制度安排」的命令、调度、照会或呈报形成独立的中央—地方行动文书链。核心取证：《清高宗实录》乾隆十六年相关条；军机处档、理藩院档或相关方略。该项锚定的是命令或决策环节，命令抵达和结果仍须另外核验。来源保留：奏折、上谕、部议、军机档或条约可证明行政动作和机构立场，不自动证明所有地区、群体或对手按其意图行动。
\eventanchor{EQ-1751-TIBET-ORDINANCE-AFTERMATH}{「清廷在拉萨事变后重整驻藏大…}乾隆十六年（1751年），「清廷在拉萨事变后重整驻藏大臣和噶厦相关制度安排」引发的善后、执行或地方回应构成可由不同材料追踪的后续节点。核心取证：《清高宗实录》乾隆十六年相关条；军机处档、理藩院档或相关方略；相关地区的档案、志书、契约、报刊或对方材料。该项锚定的是后续追踪问题，影响范围须按地区与群体分别判断。来源保留：战后、改革后或条约后的记忆、地方记录和官方总结常具有事后选择性；后续影响须按地点、群体和时间分层，不作整体化推断。
\eventanchor{EQ-1753-SALT-ADMINISTRATION-PRELUDE}{「盐政、引岸和课额问题在中央…}乾隆十八年（1753年），「盐政、引岸和课额问题在中央与产销地区持续受到整顿」之前的条件、信息和争议已通过中央与地方材料显现，构成该事件可追踪的前史节点。核心取证：《清高宗实录》乾隆十八年相关条；户部奏销、盐政、河工或海关档案；同时期地方档案、地方志、日记或对方材料应按具体地区补检。该项锚定的是前史条件与信息背景，不把条件直接写成唯一因果。来源保留：官修编年和地方材料的记录密度与立场并不相同；本行不将条件、传闻、呈报或后见解释直接等同为确定因果。
\eventanchor{EQ-1753-SALT-ADMINISTRATION-DECISION}{围绕「盐政、引岸和课额问题在…}乾隆十八年（1753年），围绕「盐政、引岸和课额问题在中央与产销地区持续受到整顿」的命令、调度、照会或呈报形成独立的中央—地方行动文书链。核心取证：《清高宗实录》乾隆十八年相关条；户部奏销、盐政、河工或海关档案。该项锚定的是命令或决策环节，命令抵达和结果仍须另外核验。来源保留：奏折、上谕、部议、军机档或条约可证明行政动作和机构立场，不自动证明所有地区、群体或对手按其意图行动。
\eventanchor{EQ-1753-SALT-ADMINISTRATION-AFTERMATH}{「盐政、引岸和课额问题在中央…}乾隆十八年（1753年），「盐政、引岸和课额问题在中央与产销地区持续受到整顿」引发的善后、执行或地方回应构成可由不同材料追踪的后续节点。核心取证：《清高宗实录》乾隆十八年相关条；户部奏销、盐政、河工或海关档案；相关地区的档案、志书、契约、报刊或对方材料。该项锚定的是后续追踪问题，影响范围须按地区与群体分别判断。来源保留：战后、改革后或条约后的记忆、地方记录和官方总结常具有事后选择性；后续影响须按地点、群体和时间分层，不作整体化推断。
\eventanchor{EQ-1754-AMURSANA-DEFECTION-PRELUDE}{「阿睦尔撒纳等准噶尔贵族向清…}乾隆十九年（1754年），「阿睦尔撒纳等准噶尔贵族向清廷归附，改变西北战争的联盟格局」之前的条件、信息和争议已通过中央与地方材料显现，构成该事件可追踪的前史节点。核心取证：《清高宗实录》乾隆十九年相关条；军机处档、理藩院档或相关方略；同时期地方档案、地方志、日记或对方材料应按具体地区补检。该项锚定的是前史条件与信息背景，不把条件直接写成唯一因果。来源保留：官修编年和地方材料的记录密度与立场并不相同；本行不将条件、传闻、呈报或后见解释直接等同为确定因果。
\eventanchor{EQ-1754-AMURSANA-DEFECTION-DECISION}{围绕「阿睦尔撒纳等准噶尔贵族…}乾隆十九年（1754年），围绕「阿睦尔撒纳等准噶尔贵族向清廷归附，改变西北战争的联盟格局」的命令、调度、照会或呈报形成独立的中央—地方行动文书链。核心取证：《清高宗实录》乾隆十九年相关条；军机处档、理藩院档或相关方略。该项锚定的是命令或决策环节，命令抵达和结果仍须另外核验。来源保留：奏折、上谕、部议、军机档或条约可证明行政动作和机构立场，不自动证明所有地区、群体或对手按其意图行动。
\eventanchor{EQ-1754-AMURSANA-DEFECTION-AFTERMATH}{「阿睦尔撒纳等准噶尔贵族向清…}乾隆十九年（1754年），「阿睦尔撒纳等准噶尔贵族向清廷归附，改变西北战争的联盟格局」引发的善后、执行或地方回应构成可由不同材料追踪的后续节点。核心取证：《清高宗实录》乾隆十九年相关条；军机处档、理藩院档或相关方略；相关地区的档案、志书、契约、报刊或对方材料。该项锚定的是后续追踪问题，影响范围须按地区与群体分别判断。来源保留：战后、改革后或条约后的记忆、地方记录和官方总结常具有事后选择性；后续影响须按地点、群体和时间分层，不作整体化推断。
\eventanchor{EQ-1755-ILI-EXPEDITION-PRELUDE}{「清军分路进入伊犁，击溃达瓦…}乾隆二十年（1755年），「清军分路进入伊犁，击溃达瓦齐政权并占领准噶尔核心地区」之前的条件、信息和争议已通过中央与地方材料显现，构成该事件可追踪的前史节点。核心取证：《清高宗实录》乾隆二十年相关条；军机处档、战事方略及地方奏报；同时期地方档案、地方志、日记或对方材料应按具体地区补检。该项锚定的是前史条件与信息背景，不把条件直接写成唯一因果。来源保留：官修编年和地方材料的记录密度与立场并不相同；本行不将条件、传闻、呈报或后见解释直接等同为确定因果。
\eventanchor{EQ-1755-ILI-EXPEDITION-DECISION}{围绕「清军分路进入伊犁，击溃…}乾隆二十年（1755年），围绕「清军分路进入伊犁，击溃达瓦齐政权并占领准噶尔核心地区」的命令、调度、照会或呈报形成独立的中央—地方行动文书链。核心取证：《清高宗实录》乾隆二十年相关条；军机处档、战事方略及地方奏报。该项锚定的是命令或决策环节，命令抵达和结果仍须另外核验。来源保留：奏折、上谕、部议、军机档或条约可证明行政动作和机构立场，不自动证明所有地区、群体或对手按其意图行动。
\eventanchor{EQ-1755-ILI-EXPEDITION-AFTERMATH}{「清军分路进入伊犁，击溃达瓦…}乾隆二十年（1755年），「清军分路进入伊犁，击溃达瓦齐政权并占领准噶尔核心地区」引发的善后、执行或地方回应构成可由不同材料追踪的后续节点。核心取证：《清高宗实录》乾隆二十年相关条；军机处档、战事方略及地方奏报；相关地区的档案、志书、契约、报刊或对方材料。该项锚定的是后续追踪问题，影响范围须按地区与群体分别判断。来源保留：战后、改革后或条约后的记忆、地方记录和官方总结常具有事后选择性；后续影响须按地点、群体和时间分层，不作整体化推断。
\eventanchor{EQ-1756-BANNER-CIVILIAN-REGISTRATION-PRELUDE}{「清廷命令部分八旗副户口改归…}乾隆二十一年（1756年），「清廷命令部分八旗副户口改归民籍，继续缓解旗籍供养压力」之前的条件、信息和争议已通过中央与地方材料显现，构成该事件可追踪的前史节点。核心取证：《清高宗实录》乾隆二十一年相关条；地方志、督抚奏折及地方档案；同时期地方档案、地方志、日记或对方材料应按具体地区补检。该项锚定的是前史条件与信息背景，不把条件直接写成唯一因果。来源保留：官修编年和地方材料的记录密度与立场并不相同；本行不将条件、传闻、呈报或后见解释直接等同为确定因果。
\eventanchor{EQ-1756-BANNER-CIVILIAN-REGISTRATION-DECISION}{围绕「清廷命令部分八旗副户口…}乾隆二十一年（1756年），围绕「清廷命令部分八旗副户口改归民籍，继续缓解旗籍供养压力」的命令、调度、照会或呈报形成独立的中央—地方行动文书链。核心取证：《清高宗实录》乾隆二十一年相关条；地方志、督抚奏折及地方档案。该项锚定的是命令或决策环节，命令抵达和结果仍须另外核验。来源保留：奏折、上谕、部议、军机档或条约可证明行政动作和机构立场，不自动证明所有地区、群体或对手按其意图行动。
\eventanchor{EQ-1756-BANNER-CIVILIAN-REGISTRATION-AFTERMATH}{「清廷命令部分八旗副户口改归…}乾隆二十一年（1756年），「清廷命令部分八旗副户口改归民籍，继续缓解旗籍供养压力」引发的善后、执行或地方回应构成可由不同材料追踪的后续节点。核心取证：《清高宗实录》乾隆二十一年相关条；地方志、督抚奏折及地方档案；相关地区的档案、志书、契约、报刊或对方材料。该项锚定的是后续追踪问题，影响范围须按地区与群体分别判断。来源保留：战后、改革后或条约后的记忆、地方记录和官方总结常具有事后选择性；后续影响须按地点、群体和时间分层，不作整体化推断。
\eventanchor{EQ-1757-CANTON-TRADE-RESTRICTION-PRELUDE}{「清廷将西洋贸易集中于广州口…}乾隆二十二年（1757年），「清廷将西洋贸易集中于广州口岸，形成后来所谓“一口通商”的制度框架」之前的条件、信息和争议已通过中央与地方材料显现，构成该事件可追踪的前史节点。核心取证：《清高宗实录》乾隆二十二年相关条；《筹办夷务始末》相关年卷与中外照会、条约文本；同时期地方档案、地方志、日记或对方材料应按具体地区补检。该项锚定的是前史条件与信息背景，不把条件直接写成唯一因果。来源保留：官修编年和地方材料的记录密度与立场并不相同；本行不将条件、传闻、呈报或后见解释直接等同为确定因果。
\eventanchor{EQ-1757-CANTON-TRADE-RESTRICTION-DECISION}{围绕「清廷将西洋贸易集中于广…}乾隆二十二年（1757年），围绕「清廷将西洋贸易集中于广州口岸，形成后来所谓“一口通商”的制度框架」的命令、调度、照会或呈报形成独立的中央—地方行动文书链。核心取证：《清高宗实录》乾隆二十二年相关条；《筹办夷务始末》相关年卷与中外照会、条约文本。该项锚定的是命令或决策环节，命令抵达和结果仍须另外核验。来源保留：奏折、上谕、部议、军机档或条约可证明行政动作和机构立场，不自动证明所有地区、群体或对手按其意图行动。
\eventanchor{EQ-1757-CANTON-TRADE-RESTRICTION-AFTERMATH}{「清廷将西洋贸易集中于广州口…}乾隆二十二年（1757年），「清廷将西洋贸易集中于广州口岸，形成后来所谓“一口通商”的制度框架」引发的善后、执行或地方回应构成可由不同材料追踪的后续节点。核心取证：《清高宗实录》乾隆二十二年相关条；《筹办夷务始末》相关年卷与中外照会、条约文本；相关地区的档案、志书、契约、报刊或对方材料。该项锚定的是后续追踪问题，影响范围须按地区与群体分别判断。来源保留：战后、改革后或条约后的记忆、地方记录和官方总结常具有事后选择性；后续影响须按地点、群体和时间分层，不作整体化推断。
\eventanchor{EQ-1757-AMURSANA-DEATH-PRELUDE}{「阿睦尔撒纳逃入俄境后死亡，…}乾隆二十二年（1757年），「阿睦尔撒纳逃入俄境后死亡，清准战争的关键对手退出」之前的条件、信息和争议已通过中央与地方材料显现，构成该事件可追踪的前史节点。核心取证：《清高宗实录》乾隆二十二年相关条；军机处档、理藩院档或相关方略；同时期地方档案、地方志、日记或对方材料应按具体地区补检。该项锚定的是前史条件与信息背景，不把条件直接写成唯一因果。来源保留：官修编年和地方材料的记录密度与立场并不相同；本行不将条件、传闻、呈报或后见解释直接等同为确定因果。
\eventanchor{EQ-1757-AMURSANA-DEATH-DECISION}{围绕「阿睦尔撒纳逃入俄境后死…}乾隆二十二年（1757年），围绕「阿睦尔撒纳逃入俄境后死亡，清准战争的关键对手退出」的命令、调度、照会或呈报形成独立的中央—地方行动文书链。核心取证：《清高宗实录》乾隆二十二年相关条；军机处档、理藩院档或相关方略。该项锚定的是命令或决策环节，命令抵达和结果仍须另外核验。来源保留：奏折、上谕、部议、军机档或条约可证明行政动作和机构立场，不自动证明所有地区、群体或对手按其意图行动。
\eventanchor{EQ-1757-AMURSANA-DEATH-AFTERMATH}{「阿睦尔撒纳逃入俄境后死亡，…}乾隆二十二年（1757年），「阿睦尔撒纳逃入俄境后死亡，清准战争的关键对手退出」引发的善后、执行或地方回应构成可由不同材料追踪的后续节点。核心取证：《清高宗实录》乾隆二十二年相关条；军机处档、理藩院档或相关方略；相关地区的档案、志书、契约、报刊或对方材料。该项锚定的是后续追踪问题，影响范围须按地区与群体分别判断。来源保留：战后、改革后或条约后的记忆、地方记录和官方总结常具有事后选择性；后续影响须按地点、群体和时间分层，不作整体化推断。
\eventanchor{EQ-1759-KASHGAR-YARKAND-CAMPAIGN-PRELUDE}{「清军攻入喀什噶尔、叶尔羌等…}乾隆二十四年（1759年），「清军攻入喀什噶尔、叶尔羌等地，和卓政权覆亡」之前的条件、信息和争议已通过中央与地方材料显现，构成该事件可追踪的前史节点。核心取证：《清高宗实录》乾隆二十四年相关条；军机处档、战事方略及地方奏报；同时期地方档案、地方志、日记或对方材料应按具体地区补检。该项锚定的是前史条件与信息背景，不把条件直接写成唯一因果。来源保留：官修编年和地方材料的记录密度与立场并不相同；本行不将条件、传闻、呈报或后见解释直接等同为确定因果。
\eventanchor{EQ-1759-KASHGAR-YARKAND-CAMPAIGN-DECISION}{围绕「清军攻入喀什噶尔、叶尔…}乾隆二十四年（1759年），围绕「清军攻入喀什噶尔、叶尔羌等地，和卓政权覆亡」的命令、调度、照会或呈报形成独立的中央—地方行动文书链。核心取证：《清高宗实录》乾隆二十四年相关条；军机处档、战事方略及地方奏报。该项锚定的是命令或决策环节，命令抵达和结果仍须另外核验。来源保留：奏折、上谕、部议、军机档或条约可证明行政动作和机构立场，不自动证明所有地区、群体或对手按其意图行动。
\eventanchor{EQ-1759-KASHGAR-YARKAND-CAMPAIGN-AFTERMATH}{「清军攻入喀什噶尔、叶尔羌等…}乾隆二十四年（1759年），「清军攻入喀什噶尔、叶尔羌等地，和卓政权覆亡」引发的善后、执行或地方回应构成可由不同材料追踪的后续节点。核心取证：《清高宗实录》乾隆二十四年相关条；军机处档、战事方略及地方奏报；相关地区的档案、志书、契约、报刊或对方材料。该项锚定的是后续追踪问题，影响范围须按地区与群体分别判断。来源保留：战后、改革后或条约后的记忆、地方记录和官方总结常具有事后选择性；后续影响须按地点、群体和时间分层，不作整体化推断。
\subsection{1760—1769}
\eventanchor{EQ-1760-ILI-GENERAL-PRELUDE}{「清廷设伊犁将军并展开驻防、…}乾隆二十五年（1760年），「清廷设伊犁将军并展开驻防、屯田与移民安置」之前的条件、信息和争议已通过中央与地方材料显现，构成该事件可追踪的前史节点。核心取证：《清高宗实录》乾隆二十五年相关条；军机处档、理藩院档或相关方略；同时期地方档案、地方志、日记或对方材料应按具体地区补检。该项锚定的是前史条件与信息背景，不把条件直接写成唯一因果。来源保留：官修编年和地方材料的记录密度与立场并不相同；本行不将条件、传闻、呈报或后见解释直接等同为确定因果。
\eventanchor{EQ-1760-ILI-GENERAL-DECISION}{围绕「清廷设伊犁将军并展开驻…}乾隆二十五年（1760年），围绕「清廷设伊犁将军并展开驻防、屯田与移民安置」的命令、调度、照会或呈报形成独立的中央—地方行动文书链。核心取证：《清高宗实录》乾隆二十五年相关条；军机处档、理藩院档或相关方略。该项锚定的是命令或决策环节，命令抵达和结果仍须另外核验。来源保留：奏折、上谕、部议、军机档或条约可证明行政动作和机构立场，不自动证明所有地区、群体或对手按其意图行动。
\eventanchor{EQ-1760-ILI-GENERAL-AFTERMATH}{「清廷设伊犁将军并展开驻防、…}乾隆二十五年（1760年），「清廷设伊犁将军并展开驻防、屯田与移民安置」引发的善后、执行或地方回应构成可由不同材料追踪的后续节点。核心取证：《清高宗实录》乾隆二十五年相关条；军机处档、理藩院档或相关方略；相关地区的档案、志书、契约、报刊或对方材料。该项锚定的是后续追踪问题，影响范围须按地区与群体分别判断。来源保留：战后、改革后或条约后的记忆、地方记录和官方总结常具有事后选择性；后续影响须按地点、群体和时间分层，不作整体化推断。
\eventanchor{EQ-1761-XINJIANG-GARRISONS-PRELUDE}{「清廷在新疆配置满汉蒙军队和…}乾隆二十六年（1761年），「清廷在新疆配置满汉蒙军队和绿营驻防，建立军政节点」之前的条件、信息和争议已通过中央与地方材料显现，构成该事件可追踪的前史节点。核心取证：《清高宗实录》乾隆二十六年相关条；军机处档、理藩院档或相关方略；同时期地方档案、地方志、日记或对方材料应按具体地区补检。该项锚定的是前史条件与信息背景，不把条件直接写成唯一因果。来源保留：官修编年和地方材料的记录密度与立场并不相同；本行不将条件、传闻、呈报或后见解释直接等同为确定因果。
\eventanchor{EQ-1761-XINJIANG-GARRISONS-DECISION}{围绕「清廷在新疆配置满汉蒙军…}乾隆二十六年（1761年），围绕「清廷在新疆配置满汉蒙军队和绿营驻防，建立军政节点」的命令、调度、照会或呈报形成独立的中央—地方行动文书链。核心取证：《清高宗实录》乾隆二十六年相关条；军机处档、理藩院档或相关方略。该项锚定的是命令或决策环节，命令抵达和结果仍须另外核验。来源保留：奏折、上谕、部议、军机档或条约可证明行政动作和机构立场，不自动证明所有地区、群体或对手按其意图行动。
\eventanchor{EQ-1761-XINJIANG-GARRISONS-AFTERMATH}{「清廷在新疆配置满汉蒙军队和…}乾隆二十六年（1761年），「清廷在新疆配置满汉蒙军队和绿营驻防，建立军政节点」引发的善后、执行或地方回应构成可由不同材料追踪的后续节点。核心取证：《清高宗实录》乾隆二十六年相关条；军机处档、理藩院档或相关方略；相关地区的档案、志书、契约、报刊或对方材料。该项锚定的是后续追踪问题，影响范围须按地区与群体分别判断。来源保留：战后、改革后或条约后的记忆、地方记录和官方总结常具有事后选择性；后续影响须按地点、群体和时间分层，不作整体化推断。
\eventanchor{EQ-1763-XINJIANG-SETTLEMENT-PRELUDE}{「清廷继续在新疆安置屯田人口…}乾隆二十八年（1763年），「清廷继续在新疆安置屯田人口、军户和商业网络」之前的条件、信息和争议已通过中央与地方材料显现，构成该事件可追踪的前史节点。核心取证：《清高宗实录》乾隆二十八年相关条；地方志、督抚奏折及地方档案；同时期地方档案、地方志、日记或对方材料应按具体地区补检。该项锚定的是前史条件与信息背景，不把条件直接写成唯一因果。来源保留：官修编年和地方材料的记录密度与立场并不相同；本行不将条件、传闻、呈报或后见解释直接等同为确定因果。
\eventanchor{EQ-1763-XINJIANG-SETTLEMENT-DECISION}{围绕「清廷继续在新疆安置屯田…}乾隆二十八年（1763年），围绕「清廷继续在新疆安置屯田人口、军户和商业网络」的命令、调度、照会或呈报形成独立的中央—地方行动文书链。核心取证：《清高宗实录》乾隆二十八年相关条；地方志、督抚奏折及地方档案。该项锚定的是命令或决策环节，命令抵达和结果仍须另外核验。来源保留：奏折、上谕、部议、军机档或条约可证明行政动作和机构立场，不自动证明所有地区、群体或对手按其意图行动。
\eventanchor{EQ-1763-XINJIANG-SETTLEMENT-AFTERMATH}{「清廷继续在新疆安置屯田人口…}乾隆二十八年（1763年），「清廷继续在新疆安置屯田人口、军户和商业网络」引发的善后、执行或地方回应构成可由不同材料追踪的后续节点。核心取证：《清高宗实录》乾隆二十八年相关条；地方志、督抚奏折及地方档案；相关地区的档案、志书、契约、报刊或对方材料。该项锚定的是后续追踪问题，影响范围须按地区与群体分别判断。来源保留：战后、改革后或条约后的记忆、地方记录和官方总结常具有事后选择性；后续影响须按地点、群体和时间分层，不作整体化推断。
\eventanchor{EQ-1764-MANCHU-BANNER-RESETTLEMENT-PRELUDE}{「部分八旗兵丁被调往西北驻防…}乾隆二十九年（1764年），「部分八旗兵丁被调往西北驻防，旗籍家庭的迁移与供给成为边务问题」之前的条件、信息和争议已通过中央与地方材料显现，构成该事件可追踪的前史节点。核心取证：《清高宗实录》乾隆二十九年相关条；地方志、督抚奏折及地方档案；同时期地方档案、地方志、日记或对方材料应按具体地区补检。该项锚定的是前史条件与信息背景，不把条件直接写成唯一因果。来源保留：官修编年和地方材料的记录密度与立场并不相同；本行不将条件、传闻、呈报或后见解释直接等同为确定因果。
\eventanchor{EQ-1764-MANCHU-BANNER-RESETTLEMENT-DECISION}{围绕「部分八旗兵丁被调往西北…}乾隆二十九年（1764年），围绕「部分八旗兵丁被调往西北驻防，旗籍家庭的迁移与供给成为边务问题」的命令、调度、照会或呈报形成独立的中央—地方行动文书链。核心取证：《清高宗实录》乾隆二十九年相关条；地方志、督抚奏折及地方档案。该项锚定的是命令或决策环节，命令抵达和结果仍须另外核验。来源保留：奏折、上谕、部议、军机档或条约可证明行政动作和机构立场，不自动证明所有地区、群体或对手按其意图行动。
\eventanchor{EQ-1764-MANCHU-BANNER-RESETTLEMENT-AFTERMATH}{「部分八旗兵丁被调往西北驻防…}乾隆二十九年（1764年），「部分八旗兵丁被调往西北驻防，旗籍家庭的迁移与供给成为边务问题」引发的善后、执行或地方回应构成可由不同材料追踪的后续节点。核心取证：《清高宗实录》乾隆二十九年相关条；地方志、督抚奏折及地方档案；相关地区的档案、志书、契约、报刊或对方材料。该项锚定的是后续追踪问题，影响范围须按地区与群体分别判断。来源保留：战后、改革后或条约后的记忆、地方记录和官方总结常具有事后选择性；后续影响须按地点、群体和时间分层，不作整体化推断。
\eventanchor{EQ-1765-USH-REBELLION-PRELUDE}{「乌什爆发反清事件，清军镇压…}乾隆三十年（1765年），「乌什爆发反清事件，清军镇压并重组当地治理」之前的条件、信息和争议已通过中央与地方材料显现，构成该事件可追踪的前史节点。核心取证：《清高宗实录》乾隆三十年相关条；军机处档、战事方略及地方奏报；同时期地方档案、地方志、日记或对方材料应按具体地区补检。该项锚定的是前史条件与信息背景，不把条件直接写成唯一因果。来源保留：官修编年和地方材料的记录密度与立场并不相同；本行不将条件、传闻、呈报或后见解释直接等同为确定因果。
\eventanchor{EQ-1765-USH-REBELLION-DECISION}{围绕「乌什爆发反清事件，清军…}乾隆三十年（1765年），围绕「乌什爆发反清事件，清军镇压并重组当地治理」的命令、调度、照会或呈报形成独立的中央—地方行动文书链。核心取证：《清高宗实录》乾隆三十年相关条；军机处档、战事方略及地方奏报。该项锚定的是命令或决策环节，命令抵达和结果仍须另外核验。来源保留：奏折、上谕、部议、军机档或条约可证明行政动作和机构立场，不自动证明所有地区、群体或对手按其意图行动。
\eventanchor{EQ-1765-USH-REBELLION-AFTERMATH}{「乌什爆发反清事件，清军镇压…}乾隆三十年（1765年），「乌什爆发反清事件，清军镇压并重组当地治理」引发的善后、执行或地方回应构成可由不同材料追踪的后续节点。核心取证：《清高宗实录》乾隆三十年相关条；军机处档、战事方略及地方奏报；相关地区的档案、志书、契约、报刊或对方材料。该项锚定的是后续追踪问题，影响范围须按地区与群体分别判断。来源保留：战后、改革后或条约后的记忆、地方记录和官方总结常具有事后选择性；后续影响须按地点、群体和时间分层，不作整体化推断。
\eventanchor{EQ-1767-BURMA-CAMPAIGN-PRELUDE}{「清缅战争持续，清军在边境和…}乾隆三十二年（1767年），「清缅战争持续，清军在边境和山地路线中遭遇重大困难」之前的条件、信息和争议已通过中央与地方材料显现，构成该事件可追踪的前史节点。核心取证：《清高宗实录》乾隆三十二年相关条；军机处档、战事方略及地方奏报；同时期地方档案、地方志、日记或对方材料应按具体地区补检。该项锚定的是前史条件与信息背景，不把条件直接写成唯一因果。来源保留：官修编年和地方材料的记录密度与立场并不相同；本行不将条件、传闻、呈报或后见解释直接等同为确定因果。
\eventanchor{EQ-1767-BURMA-CAMPAIGN-DECISION}{围绕「清缅战争持续，清军在边…}乾隆三十二年（1767年），围绕「清缅战争持续，清军在边境和山地路线中遭遇重大困难」的命令、调度、照会或呈报形成独立的中央—地方行动文书链。核心取证：《清高宗实录》乾隆三十二年相关条；军机处档、战事方略及地方奏报。该项锚定的是命令或决策环节，命令抵达和结果仍须另外核验。来源保留：奏折、上谕、部议、军机档或条约可证明行政动作和机构立场，不自动证明所有地区、群体或对手按其意图行动。
\eventanchor{EQ-1767-BURMA-CAMPAIGN-AFTERMATH}{「清缅战争持续，清军在边境和…}乾隆三十二年（1767年），「清缅战争持续，清军在边境和山地路线中遭遇重大困难」引发的善后、执行或地方回应构成可由不同材料追踪的后续节点。核心取证：《清高宗实录》乾隆三十二年相关条；军机处档、战事方略及地方奏报；相关地区的档案、志书、契约、报刊或对方材料。该项锚定的是后续追踪问题，影响范围须按地区与群体分别判断。来源保留：战后、改革后或条约后的记忆、地方记录和官方总结常具有事后选择性；后续影响须按地点、群体和时间分层，不作整体化推断。
\eventanchor{EQ-1768-LITERARY-INQUISITION-PRELUDE}{「乾隆朝以文字狱案件强化对著…}乾隆三十三年（1768年），「乾隆朝以文字狱案件强化对著述、地方士人和政治语言的控制」之前的条件、信息和争议已通过中央与地方材料显现，构成该事件可追踪的前史节点。核心取证：《清高宗实录》乾隆三十三年相关条；内阁档、文集或报刊相关记录；同时期地方档案、地方志、日记或对方材料应按具体地区补检。该项锚定的是前史条件与信息背景，不把条件直接写成唯一因果。来源保留：官修编年和地方材料的记录密度与立场并不相同；本行不将条件、传闻、呈报或后见解释直接等同为确定因果。
\eventanchor{EQ-1768-LITERARY-INQUISITION-DECISION}{围绕「乾隆朝以文字狱案件强化…}乾隆三十三年（1768年），围绕「乾隆朝以文字狱案件强化对著述、地方士人和政治语言的控制」的命令、调度、照会或呈报形成独立的中央—地方行动文书链。核心取证：《清高宗实录》乾隆三十三年相关条；内阁档、文集或报刊相关记录。该项锚定的是命令或决策环节，命令抵达和结果仍须另外核验。来源保留：奏折、上谕、部议、军机档或条约可证明行政动作和机构立场，不自动证明所有地区、群体或对手按其意图行动。
\eventanchor{EQ-1768-LITERARY-INQUISITION-AFTERMATH}{「乾隆朝以文字狱案件强化对著…}乾隆三十三年（1768年），「乾隆朝以文字狱案件强化对著述、地方士人和政治语言的控制」引发的善后、执行或地方回应构成可由不同材料追踪的后续节点。核心取证：《清高宗实录》乾隆三十三年相关条；内阁档、文集或报刊相关记录；相关地区的档案、志书、契约、报刊或对方材料。该项锚定的是后续追踪问题，影响范围须按地区与群体分别判断。来源保留：战后、改革后或条约后的记忆、地方记录和官方总结常具有事后选择性；后续影响须按地点、群体和时间分层，不作整体化推断。
\eventanchor{EQ-1768-SECOND-JINCHUAN-WAR-PRELUDE}{「第二次金川战争开始，清廷投…}乾隆三十三年（1768年），「第二次金川战争开始，清廷投入更大兵力和资源」之前的条件、信息和争议已通过中央与地方材料显现，构成该事件可追踪的前史节点。核心取证：《清高宗实录》乾隆三十三年相关条；军机处档、战事方略及地方奏报；同时期地方档案、地方志、日记或对方材料应按具体地区补检。该项锚定的是前史条件与信息背景，不把条件直接写成唯一因果。来源保留：官修编年和地方材料的记录密度与立场并不相同；本行不将条件、传闻、呈报或后见解释直接等同为确定因果。
\eventanchor{EQ-1768-SECOND-JINCHUAN-WAR-DECISION}{围绕「第二次金川战争开始，清…}乾隆三十三年（1768年），围绕「第二次金川战争开始，清廷投入更大兵力和资源」的命令、调度、照会或呈报形成独立的中央—地方行动文书链。核心取证：《清高宗实录》乾隆三十三年相关条；军机处档、战事方略及地方奏报。该项锚定的是命令或决策环节，命令抵达和结果仍须另外核验。来源保留：奏折、上谕、部议、军机档或条约可证明行政动作和机构立场，不自动证明所有地区、群体或对手按其意图行动。
\eventanchor{EQ-1768-SECOND-JINCHUAN-WAR-AFTERMATH}{「第二次金川战争开始，清廷投…}乾隆三十三年（1768年），「第二次金川战争开始，清廷投入更大兵力和资源」引发的善后、执行或地方回应构成可由不同材料追踪的后续节点。核心取证：《清高宗实录》乾隆三十三年相关条；军机处档、战事方略及地方奏报；相关地区的档案、志书、契约、报刊或对方材料。该项锚定的是后续追踪问题，影响范围须按地区与群体分别判断。来源保留：战后、改革后或条约后的记忆、地方记录和官方总结常具有事后选择性；后续影响须按地点、群体和时间分层，不作整体化推断。
\eventanchor{EQ-1769-JINCHUAN-ESCALATION-PRELUDE}{「金川战场的道路、炮兵、粮运…}乾隆三十四年（1769年），「金川战场的道路、炮兵、粮运与将领调度不断升级」之前的条件、信息和争议已通过中央与地方材料显现，构成该事件可追踪的前史节点。核心取证：《清高宗实录》乾隆三十四年相关条；军机处档、战事方略及地方奏报；同时期地方档案、地方志、日记或对方材料应按具体地区补检。该项锚定的是前史条件与信息背景，不把条件直接写成唯一因果。来源保留：官修编年和地方材料的记录密度与立场并不相同；本行不将条件、传闻、呈报或后见解释直接等同为确定因果。
\eventanchor{EQ-1769-JINCHUAN-ESCALATION-DECISION}{围绕「金川战场的道路、炮兵、…}乾隆三十四年（1769年），围绕「金川战场的道路、炮兵、粮运与将领调度不断升级」的命令、调度、照会或呈报形成独立的中央—地方行动文书链。核心取证：《清高宗实录》乾隆三十四年相关条；军机处档、战事方略及地方奏报。该项锚定的是命令或决策环节，命令抵达和结果仍须另外核验。来源保留：奏折、上谕、部议、军机档或条约可证明行政动作和机构立场，不自动证明所有地区、群体或对手按其意图行动。
\eventanchor{EQ-1769-JINCHUAN-ESCALATION-AFTERMATH}{「金川战场的道路、炮兵、粮运…}乾隆三十四年（1769年），「金川战场的道路、炮兵、粮运与将领调度不断升级」引发的善后、执行或地方回应构成可由不同材料追踪的后续节点。核心取证：《清高宗实录》乾隆三十四年相关条；军机处档、战事方略及地方奏报；相关地区的档案、志书、契约、报刊或对方材料。该项锚定的是后续追踪问题，影响范围须按地区与群体分别判断。来源保留：战后、改革后或条约后的记忆、地方记录和官方总结常具有事后选择性；后续影响须按地点、群体和时间分层，不作整体化推断。
\subsection{1770—1779}
\eventanchor{EQ-1770-SICHUAN-SUPPLY-BURDEN-PRELUDE}{「金川军需加重四川及邻省的转…}乾隆三十五年（1770年），「金川军需加重四川及邻省的转运、民夫和财政负担」之前的条件、信息和争议已通过中央与地方材料显现，构成该事件可追踪的前史节点。核心取证：《清高宗实录》乾隆三十五年相关条；户部奏销、盐政、河工或海关档案；同时期地方档案、地方志、日记或对方材料应按具体地区补检。该项锚定的是前史条件与信息背景，不把条件直接写成唯一因果。来源保留：官修编年和地方材料的记录密度与立场并不相同；本行不将条件、传闻、呈报或后见解释直接等同为确定因果。
\eventanchor{EQ-1770-SICHUAN-SUPPLY-BURDEN-DECISION}{围绕「金川军需加重四川及邻省…}乾隆三十五年（1770年），围绕「金川军需加重四川及邻省的转运、民夫和财政负担」的命令、调度、照会或呈报形成独立的中央—地方行动文书链。核心取证：《清高宗实录》乾隆三十五年相关条；户部奏销、盐政、河工或海关档案。该项锚定的是命令或决策环节，命令抵达和结果仍须另外核验。来源保留：奏折、上谕、部议、军机档或条约可证明行政动作和机构立场，不自动证明所有地区、群体或对手按其意图行动。
\eventanchor{EQ-1770-SICHUAN-SUPPLY-BURDEN-AFTERMATH}{「金川军需加重四川及邻省的转…}乾隆三十五年（1770年），「金川军需加重四川及邻省的转运、民夫和财政负担」引发的善后、执行或地方回应构成可由不同材料追踪的后续节点。核心取证：《清高宗实录》乾隆三十五年相关条；户部奏销、盐政、河工或海关档案；相关地区的档案、志书、契约、报刊或对方材料。该项锚定的是后续追踪问题，影响范围须按地区与群体分别判断。来源保留：战后、改革后或条约后的记忆、地方记录和官方总结常具有事后选择性；后续影响须按地点、群体和时间分层，不作整体化推断。
\eventanchor{EQ-1771-TORGHUT-RETURN-PRELUDE}{「土尔扈特部自伏尔加回归，清…}乾隆三十六年（1771年），「土尔扈特部自伏尔加回归，清廷在伊犁等地安置部众」之前的条件、信息和争议已通过中央与地方材料显现，构成该事件可追踪的前史节点。核心取证：《清高宗实录》乾隆三十六年相关条；军机处档、理藩院档或相关方略；同时期地方档案、地方志、日记或对方材料应按具体地区补检。该项锚定的是前史条件与信息背景，不把条件直接写成唯一因果。来源保留：官修编年和地方材料的记录密度与立场并不相同；本行不将条件、传闻、呈报或后见解释直接等同为确定因果。
\eventanchor{EQ-1771-TORGHUT-RETURN-DECISION}{围绕「土尔扈特部自伏尔加回归…}乾隆三十六年（1771年），围绕「土尔扈特部自伏尔加回归，清廷在伊犁等地安置部众」的命令、调度、照会或呈报形成独立的中央—地方行动文书链。核心取证：《清高宗实录》乾隆三十六年相关条；军机处档、理藩院档或相关方略。该项锚定的是命令或决策环节，命令抵达和结果仍须另外核验。来源保留：奏折、上谕、部议、军机档或条约可证明行政动作和机构立场，不自动证明所有地区、群体或对手按其意图行动。
\eventanchor{EQ-1771-TORGHUT-RETURN-AFTERMATH}{「土尔扈特部自伏尔加回归，清…}乾隆三十六年（1771年），「土尔扈特部自伏尔加回归，清廷在伊犁等地安置部众」引发的善后、执行或地方回应构成可由不同材料追踪的后续节点。核心取证：《清高宗实录》乾隆三十六年相关条；军机处档、理藩院档或相关方略；相关地区的档案、志书、契约、报刊或对方材料。该项锚定的是后续追踪问题，影响范围须按地区与群体分别判断。来源保留：战后、改革后或条约后的记忆、地方记录和官方总结常具有事后选择性；后续影响须按地点、群体和时间分层，不作整体化推断。
\eventanchor{EQ-1772-SIKU-QUANSHU-ORDER-PRELUDE}{「乾隆帝下令征集、编纂《四库…}乾隆三十七年（1772年），「乾隆帝下令征集、编纂《四库全书》，建立全国性的图书征集机制」之前的条件、信息和争议已通过中央与地方材料显现，构成该事件可追踪的前史节点。核心取证：《清高宗实录》乾隆三十七年相关条；内阁档、文集或报刊相关记录；同时期地方档案、地方志、日记或对方材料应按具体地区补检。该项锚定的是前史条件与信息背景，不把条件直接写成唯一因果。来源保留：官修编年和地方材料的记录密度与立场并不相同；本行不将条件、传闻、呈报或后见解释直接等同为确定因果。
\eventanchor{EQ-1772-SIKU-QUANSHU-ORDER-DECISION}{围绕「乾隆帝下令征集、编纂《…}乾隆三十七年（1772年），围绕「乾隆帝下令征集、编纂《四库全书》，建立全国性的图书征集机制」的命令、调度、照会或呈报形成独立的中央—地方行动文书链。核心取证：《清高宗实录》乾隆三十七年相关条；内阁档、文集或报刊相关记录。该项锚定的是命令或决策环节，命令抵达和结果仍须另外核验。来源保留：奏折、上谕、部议、军机档或条约可证明行政动作和机构立场，不自动证明所有地区、群体或对手按其意图行动。
\eventanchor{EQ-1772-SIKU-QUANSHU-ORDER-AFTERMATH}{「乾隆帝下令征集、编纂《四库…}乾隆三十七年（1772年），「乾隆帝下令征集、编纂《四库全书》，建立全国性的图书征集机制」引发的善后、执行或地方回应构成可由不同材料追踪的后续节点。核心取证：《清高宗实录》乾隆三十七年相关条；内阁档、文集或报刊相关记录；相关地区的档案、志书、契约、报刊或对方材料。该项锚定的是后续追踪问题，影响范围须按地区与群体分别判断。来源保留：战后、改革后或条约后的记忆、地方记录和官方总结常具有事后选择性；后续影响须按地点、群体和时间分层，不作整体化推断。
\eventanchor{EQ-1774-WANG-LUN-UPRISING-PRELUDE}{「王伦在山东起事，清廷迅速调…}乾隆三十九年（1774年），「王伦在山东起事，清廷迅速调兵镇压」之前的条件、信息和争议已通过中央与地方材料显现，构成该事件可追踪的前史节点。核心取证：《清高宗实录》乾隆三十九年相关条；地方志、督抚奏折及地方档案；同时期地方档案、地方志、日记或对方材料应按具体地区补检。该项锚定的是前史条件与信息背景，不把条件直接写成唯一因果。来源保留：官修编年和地方材料的记录密度与立场并不相同；本行不将条件、传闻、呈报或后见解释直接等同为确定因果。
\eventanchor{EQ-1774-WANG-LUN-UPRISING-DECISION}{围绕「王伦在山东起事，清廷迅…}乾隆三十九年（1774年），围绕「王伦在山东起事，清廷迅速调兵镇压」的命令、调度、照会或呈报形成独立的中央—地方行动文书链。核心取证：《清高宗实录》乾隆三十九年相关条；地方志、督抚奏折及地方档案。该项锚定的是命令或决策环节，命令抵达和结果仍须另外核验。来源保留：奏折、上谕、部议、军机档或条约可证明行政动作和机构立场，不自动证明所有地区、群体或对手按其意图行动。
\eventanchor{EQ-1774-WANG-LUN-UPRISING-AFTERMATH}{「王伦在山东起事，清廷迅速调…}乾隆三十九年（1774年），「王伦在山东起事，清廷迅速调兵镇压」引发的善后、执行或地方回应构成可由不同材料追踪的后续节点。核心取证：《清高宗实录》乾隆三十九年相关条；地方志、督抚奏折及地方档案；相关地区的档案、志书、契约、报刊或对方材料。该项锚定的是后续追踪问题，影响范围须按地区与群体分别判断。来源保留：战后、改革后或条约后的记忆、地方记录和官方总结常具有事后选择性；后续影响须按地点、群体和时间分层，不作整体化推断。
\eventanchor{EQ-1775-SIKU-COLLECTION-PRELUDE}{「《四库全书》征书与审查在各…}乾隆四十年（1775年），「《四库全书》征书与审查在各省展开，地方藏书和文人网络受到影响」之前的条件、信息和争议已通过中央与地方材料显现，构成该事件可追踪的前史节点。核心取证：《清高宗实录》乾隆四十年相关条；内阁档、文集或报刊相关记录；同时期地方档案、地方志、日记或对方材料应按具体地区补检。该项锚定的是前史条件与信息背景，不把条件直接写成唯一因果。来源保留：官修编年和地方材料的记录密度与立场并不相同；本行不将条件、传闻、呈报或后见解释直接等同为确定因果。
\eventanchor{EQ-1775-SIKU-COLLECTION-DECISION}{围绕「《四库全书》征书与审查…}乾隆四十年（1775年），围绕「《四库全书》征书与审查在各省展开，地方藏书和文人网络受到影响」的命令、调度、照会或呈报形成独立的中央—地方行动文书链。核心取证：《清高宗实录》乾隆四十年相关条；内阁档、文集或报刊相关记录。该项锚定的是命令或决策环节，命令抵达和结果仍须另外核验。来源保留：奏折、上谕、部议、军机档或条约可证明行政动作和机构立场，不自动证明所有地区、群体或对手按其意图行动。
\eventanchor{EQ-1775-SIKU-COLLECTION-AFTERMATH}{「《四库全书》征书与审查在各…}乾隆四十年（1775年），「《四库全书》征书与审查在各省展开，地方藏书和文人网络受到影响」引发的善后、执行或地方回应构成可由不同材料追踪的后续节点。核心取证：《清高宗实录》乾隆四十年相关条；内阁档、文集或报刊相关记录；相关地区的档案、志书、契约、报刊或对方材料。该项锚定的是后续追踪问题，影响范围须按地区与群体分别判断。来源保留：战后、改革后或条约后的记忆、地方记录和官方总结常具有事后选择性；后续影响须按地点、群体和时间分层，不作整体化推断。
\eventanchor{EQ-1776-JINCHUAN-CONQUEST-PRELUDE}{「清军攻克金川，第二次金川战…}乾隆四十一年（1776年），「清军攻克金川，第二次金川战争结束」之前的条件、信息和争议已通过中央与地方材料显现，构成该事件可追踪的前史节点。核心取证：《清高宗实录》乾隆四十一年相关条；军机处档、战事方略及地方奏报；同时期地方档案、地方志、日记或对方材料应按具体地区补检。该项锚定的是前史条件与信息背景，不把条件直接写成唯一因果。来源保留：官修编年和地方材料的记录密度与立场并不相同；本行不将条件、传闻、呈报或后见解释直接等同为确定因果。
\eventanchor{EQ-1776-JINCHUAN-CONQUEST-DECISION}{围绕「清军攻克金川，第二次金…}乾隆四十一年（1776年），围绕「清军攻克金川，第二次金川战争结束」的命令、调度、照会或呈报形成独立的中央—地方行动文书链。核心取证：《清高宗实录》乾隆四十一年相关条；军机处档、战事方略及地方奏报。该项锚定的是命令或决策环节，命令抵达和结果仍须另外核验。来源保留：奏折、上谕、部议、军机档或条约可证明行政动作和机构立场，不自动证明所有地区、群体或对手按其意图行动。
\eventanchor{EQ-1776-JINCHUAN-CONQUEST-AFTERMATH}{「清军攻克金川，第二次金川战…}乾隆四十一年（1776年），「清军攻克金川，第二次金川战争结束」引发的善后、执行或地方回应构成可由不同材料追踪的后续节点。核心取证：《清高宗实录》乾隆四十一年相关条；军机处档、战事方略及地方奏报；相关地区的档案、志书、契约、报刊或对方材料。该项锚定的是后续追踪问题，影响范围须按地区与群体分别判断。来源保留：战后、改革后或条约后的记忆、地方记录和官方总结常具有事后选择性；后续影响须按地点、群体和时间分层，不作整体化推断。
\eventanchor{EQ-1778-SICHUAN-TAX-RECOVERY-PRELUDE}{「金川战后四川财政和民力恢复…}乾隆四十三年（1778年），「金川战后四川财政和民力恢复成为地方官奏报的重要议题」之前的条件、信息和争议已通过中央与地方材料显现，构成该事件可追踪的前史节点。核心取证：《清高宗实录》乾隆四十三年相关条；户部奏销、盐政、河工或海关档案；同时期地方档案、地方志、日记或对方材料应按具体地区补检。该项锚定的是前史条件与信息背景，不把条件直接写成唯一因果。来源保留：官修编年和地方材料的记录密度与立场并不相同；本行不将条件、传闻、呈报或后见解释直接等同为确定因果。
\eventanchor{EQ-1778-SICHUAN-TAX-RECOVERY-DECISION}{围绕「金川战后四川财政和民力…}乾隆四十三年（1778年），围绕「金川战后四川财政和民力恢复成为地方官奏报的重要议题」的命令、调度、照会或呈报形成独立的中央—地方行动文书链。核心取证：《清高宗实录》乾隆四十三年相关条；户部奏销、盐政、河工或海关档案。该项锚定的是命令或决策环节，命令抵达和结果仍须另外核验。来源保留：奏折、上谕、部议、军机档或条约可证明行政动作和机构立场，不自动证明所有地区、群体或对手按其意图行动。
\eventanchor{EQ-1778-SICHUAN-TAX-RECOVERY-AFTERMATH}{「金川战后四川财政和民力恢复…}乾隆四十三年（1778年），「金川战后四川财政和民力恢复成为地方官奏报的重要议题」引发的善后、执行或地方回应构成可由不同材料追踪的后续节点。核心取证：《清高宗实录》乾隆四十三年相关条；户部奏销、盐政、河工或海关档案；相关地区的档案、志书、契约、报刊或对方材料。该项锚定的是后续追踪问题，影响范围须按地区与群体分别判断。来源保留：战后、改革后或条约后的记忆、地方记录和官方总结常具有事后选择性；后续影响须按地点、群体和时间分层，不作整体化推断。
\eventanchor{EQ-1779-VIETNAM-INTERVENTION-PRELUDE}{「清廷介入安南政局并派兵南下…}乾隆四十四年（1779年），「清廷介入安南政局并派兵南下，孙士毅军一度进入升龙」之前的条件、信息和争议已通过中央与地方材料显现，构成该事件可追踪的前史节点。核心取证：《清高宗实录》乾隆四十四年相关条；军机处档、战事方略及地方奏报；同时期地方档案、地方志、日记或对方材料应按具体地区补检。该项锚定的是前史条件与信息背景，不把条件直接写成唯一因果。来源保留：官修编年和地方材料的记录密度与立场并不相同；本行不将条件、传闻、呈报或后见解释直接等同为确定因果。
\eventanchor{EQ-1779-VIETNAM-INTERVENTION-DECISION}{围绕「清廷介入安南政局并派兵…}乾隆四十四年（1779年），围绕「清廷介入安南政局并派兵南下，孙士毅军一度进入升龙」的命令、调度、照会或呈报形成独立的中央—地方行动文书链。核心取证：《清高宗实录》乾隆四十四年相关条；军机处档、战事方略及地方奏报。该项锚定的是命令或决策环节，命令抵达和结果仍须另外核验。来源保留：奏折、上谕、部议、军机档或条约可证明行政动作和机构立场，不自动证明所有地区、群体或对手按其意图行动。
\eventanchor{EQ-1779-VIETNAM-INTERVENTION-AFTERMATH}{「清廷介入安南政局并派兵南下…}乾隆四十四年（1779年），「清廷介入安南政局并派兵南下，孙士毅军一度进入升龙」引发的善后、执行或地方回应构成可由不同材料追踪的后续节点。核心取证：《清高宗实录》乾隆四十四年相关条；军机处档、战事方略及地方奏报；相关地区的档案、志书、契约、报刊或对方材料。该项锚定的是后续追踪问题，影响范围须按地区与群体分别判断。来源保留：战后、改革后或条约后的记忆、地方记录和官方总结常具有事后选择性；后续影响须按地点、群体和时间分层，不作整体化推断。
\subsection{1780—1789}
\eventanchor{EQ-1789-TIBET-NEPAL-CAMPAIGN-DOCUMENTS}{清廷围绕「清军越过喜马拉雅方…}乾隆五十四年（1789年），清廷围绕「清军越过喜马拉雅方向追击廓尔喀，边境协商与军事行动并行」调遣军队、筹措饷械并处理战报，中央与地方的军政文书形成连续记录。核心取证：《清高宗实录》乾隆五十四年相关条；军机处档、战事方略及地方奏报。该项锚定的是可回查的文书动作，不将其直接等同于地方执行。来源保留：此行仅确认相关文书链和可追查的行政动作；不得由其直接推断政策有效、地方服从、民众态度或单一因果。
\eventanchor{EQ-1780-VIETNAM-WITHDRAWAL-PRELUDE}{「清军撤离安南，双方通过使节…}乾隆四十五年（1780年），「清军撤离安南，双方通过使节和册封语言恢复关系」之前的条件、信息和争议已通过中央与地方材料显现，构成该事件可追踪的前史节点。核心取证：《清高宗实录》乾隆四十五年相关条；《筹办夷务始末》相关年卷与中外照会、条约文本；同时期地方档案、地方志、日记或对方材料应按具体地区补检。该项锚定的是前史条件与信息背景，不把条件直接写成唯一因果。来源保留：官修编年和地方材料的记录密度与立场并不相同；本行不将条件、传闻、呈报或后见解释直接等同为确定因果。
\eventanchor{EQ-1780-VIETNAM-WITHDRAWAL-DECISION}{围绕「清军撤离安南，双方通过…}乾隆四十五年（1780年），围绕「清军撤离安南，双方通过使节和册封语言恢复关系」的命令、调度、照会或呈报形成独立的中央—地方行动文书链。核心取证：《清高宗实录》乾隆四十五年相关条；《筹办夷务始末》相关年卷与中外照会、条约文本。该项锚定的是命令或决策环节，命令抵达和结果仍须另外核验。来源保留：奏折、上谕、部议、军机档或条约可证明行政动作和机构立场，不自动证明所有地区、群体或对手按其意图行动。
\eventanchor{EQ-1780-VIETNAM-WITHDRAWAL-AFTERMATH}{「清军撤离安南，双方通过使节…}乾隆四十五年（1780年），「清军撤离安南，双方通过使节和册封语言恢复关系」引发的善后、执行或地方回应构成可由不同材料追踪的后续节点。核心取证：《清高宗实录》乾隆四十五年相关条；《筹办夷务始末》相关年卷与中外照会、条约文本；相关地区的档案、志书、契约、报刊或对方材料。该项锚定的是后续追踪问题，影响范围须按地区与群体分别判断。来源保留：战后、改革后或条约后的记忆、地方记录和官方总结常具有事后选择性；后续影响须按地点、群体和时间分层，不作整体化推断。
\eventanchor{EQ-1782-GANSU-POSTWAR-GOVERNANCE-PRELUDE}{「清廷在甘肃起事后调整驻兵、…}乾隆四十七年（1782年），「清廷在甘肃起事后调整驻兵、司法和地方控制措施」之前的条件、信息和争议已通过中央与地方材料显现，构成该事件可追踪的前史节点。核心取证：《清高宗实录》乾隆四十七年相关条；地方志、督抚奏折及地方档案；同时期地方档案、地方志、日记或对方材料应按具体地区补检。该项锚定的是前史条件与信息背景，不把条件直接写成唯一因果。来源保留：官修编年和地方材料的记录密度与立场并不相同；本行不将条件、传闻、呈报或后见解释直接等同为确定因果。
\eventanchor{EQ-1782-GANSU-POSTWAR-GOVERNANCE-DECISION}{围绕「清廷在甘肃起事后调整驻…}乾隆四十七年（1782年），围绕「清廷在甘肃起事后调整驻兵、司法和地方控制措施」的命令、调度、照会或呈报形成独立的中央—地方行动文书链。核心取证：《清高宗实录》乾隆四十七年相关条；地方志、督抚奏折及地方档案。该项锚定的是命令或决策环节，命令抵达和结果仍须另外核验。来源保留：奏折、上谕、部议、军机档或条约可证明行政动作和机构立场，不自动证明所有地区、群体或对手按其意图行动。
\eventanchor{EQ-1782-GANSU-POSTWAR-GOVERNANCE-AFTERMATH}{「清廷在甘肃起事后调整驻兵、…}乾隆四十七年（1782年），「清廷在甘肃起事后调整驻兵、司法和地方控制措施」引发的善后、执行或地方回应构成可由不同材料追踪的后续节点。核心取证：《清高宗实录》乾隆四十七年相关条；地方志、督抚奏折及地方档案；相关地区的档案、志书、契约、报刊或对方材料。该项锚定的是后续追踪问题，影响范围须按地区与群体分别判断。来源保留：战后、改革后或条约后的记忆、地方记录和官方总结常具有事后选择性；后续影响须按地点、群体和时间分层，不作整体化推断。
\eventanchor{EQ-1783-SALT-MONOPOLY-REVIEW-PRELUDE}{「清廷继续讨论盐政弊端、商人…}乾隆四十八年（1783年），「清廷继续讨论盐政弊端、商人负担和课额征解」之前的条件、信息和争议已通过中央与地方材料显现，构成该事件可追踪的前史节点。核心取证：《清高宗实录》乾隆四十八年相关条；户部奏销、盐政、河工或海关档案；同时期地方档案、地方志、日记或对方材料应按具体地区补检。该项锚定的是前史条件与信息背景，不把条件直接写成唯一因果。来源保留：官修编年和地方材料的记录密度与立场并不相同；本行不将条件、传闻、呈报或后见解释直接等同为确定因果。
\eventanchor{EQ-1783-SALT-MONOPOLY-REVIEW-DECISION}{围绕「清廷继续讨论盐政弊端、…}乾隆四十八年（1783年），围绕「清廷继续讨论盐政弊端、商人负担和课额征解」的命令、调度、照会或呈报形成独立的中央—地方行动文书链。核心取证：《清高宗实录》乾隆四十八年相关条；户部奏销、盐政、河工或海关档案。该项锚定的是命令或决策环节，命令抵达和结果仍须另外核验。来源保留：奏折、上谕、部议、军机档或条约可证明行政动作和机构立场，不自动证明所有地区、群体或对手按其意图行动。
\eventanchor{EQ-1783-SALT-MONOPOLY-REVIEW-AFTERMATH}{「清廷继续讨论盐政弊端、商人…}乾隆四十八年（1783年），「清廷继续讨论盐政弊端、商人负担和课额征解」引发的善后、执行或地方回应构成可由不同材料追踪的后续节点。核心取证：《清高宗实录》乾隆四十八年相关条；户部奏销、盐政、河工或海关档案；相关地区的档案、志书、契约、报刊或对方材料。该项锚定的是后续追踪问题，影响范围须按地区与群体分别判断。来源保留：战后、改革后或条约后的记忆、地方记录和官方总结常具有事后选择性；后续影响须按地点、群体和时间分层，不作整体化推断。
\eventanchor{EQ-1784-TAIWAN-UPRISING-PRELUDE-PRELUDE}{「台湾地方械斗、赋役与移民社…}乾隆四十九年（1784年），「台湾地方械斗、赋役与移民社会紧张为大规模反抗累积条件」之前的条件、信息和争议已通过中央与地方材料显现，构成该事件可追踪的前史节点。核心取证：《清高宗实录》乾隆四十九年相关条；地方志、督抚奏折及地方档案；同时期地方档案、地方志、日记或对方材料应按具体地区补检。该项锚定的是前史条件与信息背景，不把条件直接写成唯一因果。来源保留：官修编年和地方材料的记录密度与立场并不相同；本行不将条件、传闻、呈报或后见解释直接等同为确定因果。
\eventanchor{EQ-1784-TAIWAN-UPRISING-PRELUDE-DECISION}{围绕「台湾地方械斗、赋役与移…}乾隆四十九年（1784年），围绕「台湾地方械斗、赋役与移民社会紧张为大规模反抗累积条件」的命令、调度、照会或呈报形成独立的中央—地方行动文书链。核心取证：《清高宗实录》乾隆四十九年相关条；地方志、督抚奏折及地方档案。该项锚定的是命令或决策环节，命令抵达和结果仍须另外核验。来源保留：奏折、上谕、部议、军机档或条约可证明行政动作和机构立场，不自动证明所有地区、群体或对手按其意图行动。
\eventanchor{EQ-1784-TAIWAN-UPRISING-PRELUDE-AFTERMATH}{「台湾地方械斗、赋役与移民社…}乾隆四十九年（1784年），「台湾地方械斗、赋役与移民社会紧张为大规模反抗累积条件」引发的善后、执行或地方回应构成可由不同材料追踪的后续节点。核心取证：《清高宗实录》乾隆四十九年相关条；地方志、督抚奏折及地方档案；相关地区的档案、志书、契约、报刊或对方材料。该项锚定的是后续追踪问题，影响范围须按地区与群体分别判断。来源保留：战后、改革后或条约后的记忆、地方记录和官方总结常具有事后选择性；后续影响须按地点、群体和时间分层，不作整体化推断。
\eventanchor{EQ-1786-LIN-SHUANGWEN-UPRISING-PRELUDE}{「林爽文起事在台湾爆发并迅速…}乾隆五十一年（1786年），「林爽文起事在台湾爆发并迅速扩展」之前的条件、信息和争议已通过中央与地方材料显现，构成该事件可追踪的前史节点。核心取证：《清高宗实录》乾隆五十一年相关条；军机处档、战事方略及地方奏报；同时期地方档案、地方志、日记或对方材料应按具体地区补检。该项锚定的是前史条件与信息背景，不把条件直接写成唯一因果。来源保留：官修编年和地方材料的记录密度与立场并不相同；本行不将条件、传闻、呈报或后见解释直接等同为确定因果。
\eventanchor{EQ-1786-LIN-SHUANGWEN-UPRISING-DECISION}{围绕「林爽文起事在台湾爆发并…}乾隆五十一年（1786年），围绕「林爽文起事在台湾爆发并迅速扩展」的命令、调度、照会或呈报形成独立的中央—地方行动文书链。核心取证：《清高宗实录》乾隆五十一年相关条；军机处档、战事方略及地方奏报。该项锚定的是命令或决策环节，命令抵达和结果仍须另外核验。来源保留：奏折、上谕、部议、军机档或条约可证明行政动作和机构立场，不自动证明所有地区、群体或对手按其意图行动。
\eventanchor{EQ-1786-LIN-SHUANGWEN-UPRISING-AFTERMATH}{「林爽文起事在台湾爆发并迅速…}乾隆五十一年（1786年），「林爽文起事在台湾爆发并迅速扩展」引发的善后、执行或地方回应构成可由不同材料追踪的后续节点。核心取证：《清高宗实录》乾隆五十一年相关条；军机处档、战事方略及地方奏报；相关地区的档案、志书、契约、报刊或对方材料。该项锚定的是后续追踪问题，影响范围须按地区与群体分别判断。来源保留：战后、改革后或条约后的记忆、地方记录和官方总结常具有事后选择性；后续影响须按地点、群体和时间分层，不作整体化推断。
\eventanchor{EQ-1786-FIRST-GURKHA-WAR-PRELUDE}{「廓尔喀军进入西藏边境，清廷…}乾隆五十一年（1786年），「廓尔喀军进入西藏边境，清廷开始处理高原军事危机」之前的条件、信息和争议已通过中央与地方材料显现，构成该事件可追踪的前史节点。核心取证：《清高宗实录》乾隆五十一年相关条；军机处档、理藩院档或相关方略；同时期地方档案、地方志、日记或对方材料应按具体地区补检。该项锚定的是前史条件与信息背景，不把条件直接写成唯一因果。来源保留：官修编年和地方材料的记录密度与立场并不相同；本行不将条件、传闻、呈报或后见解释直接等同为确定因果。
\eventanchor{EQ-1786-FIRST-GURKHA-WAR-DECISION}{围绕「廓尔喀军进入西藏边境，…}乾隆五十一年（1786年），围绕「廓尔喀军进入西藏边境，清廷开始处理高原军事危机」的命令、调度、照会或呈报形成独立的中央—地方行动文书链。核心取证：《清高宗实录》乾隆五十一年相关条；军机处档、理藩院档或相关方略。该项锚定的是命令或决策环节，命令抵达和结果仍须另外核验。来源保留：奏折、上谕、部议、军机档或条约可证明行政动作和机构立场，不自动证明所有地区、群体或对手按其意图行动。
\eventanchor{EQ-1786-FIRST-GURKHA-WAR-AFTERMATH}{「廓尔喀军进入西藏边境，清廷…}乾隆五十一年（1786年），「廓尔喀军进入西藏边境，清廷开始处理高原军事危机」引发的善后、执行或地方回应构成可由不同材料追踪的后续节点。核心取证：《清高宗实录》乾隆五十一年相关条；军机处档、理藩院档或相关方略；相关地区的档案、志书、契约、报刊或对方材料。该项锚定的是后续追踪问题，影响范围须按地区与群体分别判断。来源保留：战后、改革后或条约后的记忆、地方记录和官方总结常具有事后选择性；后续影响须按地点、群体和时间分层，不作整体化推断。
\eventanchor{EQ-1787-TAIWAN-CAMPAIGN-PRELUDE}{「福康安等率军赴台镇压林爽文…}乾隆五十二年（1787年），「福康安等率军赴台镇压林爽文起事，海峡运输成为关键」之前的条件、信息和争议已通过中央与地方材料显现，构成该事件可追踪的前史节点。核心取证：《清高宗实录》乾隆五十二年相关条；军机处档、战事方略及地方奏报；同时期地方档案、地方志、日记或对方材料应按具体地区补检。该项锚定的是前史条件与信息背景，不把条件直接写成唯一因果。来源保留：官修编年和地方材料的记录密度与立场并不相同；本行不将条件、传闻、呈报或后见解释直接等同为确定因果。
\eventanchor{EQ-1787-TAIWAN-CAMPAIGN-DECISION}{围绕「福康安等率军赴台镇压林…}乾隆五十二年（1787年），围绕「福康安等率军赴台镇压林爽文起事，海峡运输成为关键」的命令、调度、照会或呈报形成独立的中央—地方行动文书链。核心取证：《清高宗实录》乾隆五十二年相关条；军机处档、战事方略及地方奏报。该项锚定的是命令或决策环节，命令抵达和结果仍须另外核验。来源保留：奏折、上谕、部议、军机档或条约可证明行政动作和机构立场，不自动证明所有地区、群体或对手按其意图行动。
\eventanchor{EQ-1787-TAIWAN-CAMPAIGN-AFTERMATH}{「福康安等率军赴台镇压林爽文…}乾隆五十二年（1787年），「福康安等率军赴台镇压林爽文起事，海峡运输成为关键」引发的善后、执行或地方回应构成可由不同材料追踪的后续节点。核心取证：《清高宗实录》乾隆五十二年相关条；军机处档、战事方略及地方奏报；相关地区的档案、志书、契约、报刊或对方材料。该项锚定的是后续追踪问题，影响范围须按地区与群体分别判断。来源保留：战后、改革后或条约后的记忆、地方记录和官方总结常具有事后选择性；后续影响须按地点、群体和时间分层，不作整体化推断。
\eventanchor{EQ-1788-SECOND-GURKHA-WAR-PRELUDE}{「廓尔喀再次入侵西藏，清军组…}乾隆五十三年（1788年），「廓尔喀再次入侵西藏，清军组织更大规模高原远征」之前的条件、信息和争议已通过中央与地方材料显现，构成该事件可追踪的前史节点。核心取证：《清高宗实录》乾隆五十三年相关条；军机处档、战事方略及地方奏报；同时期地方档案、地方志、日记或对方材料应按具体地区补检。该项锚定的是前史条件与信息背景，不把条件直接写成唯一因果。来源保留：官修编年和地方材料的记录密度与立场并不相同；本行不将条件、传闻、呈报或后见解释直接等同为确定因果。
\eventanchor{EQ-1788-SECOND-GURKHA-WAR-DECISION}{围绕「廓尔喀再次入侵西藏，清…}乾隆五十三年（1788年），围绕「廓尔喀再次入侵西藏，清军组织更大规模高原远征」的命令、调度、照会或呈报形成独立的中央—地方行动文书链。核心取证：《清高宗实录》乾隆五十三年相关条；军机处档、战事方略及地方奏报。该项锚定的是命令或决策环节，命令抵达和结果仍须另外核验。来源保留：奏折、上谕、部议、军机档或条约可证明行政动作和机构立场，不自动证明所有地区、群体或对手按其意图行动。
\eventanchor{EQ-1788-SECOND-GURKHA-WAR-AFTERMATH}{「廓尔喀再次入侵西藏，清军组…}乾隆五十三年（1788年），「廓尔喀再次入侵西藏，清军组织更大规模高原远征」引发的善后、执行或地方回应构成可由不同材料追踪的后续节点。核心取证：《清高宗实录》乾隆五十三年相关条；军机处档、战事方略及地方奏报；相关地区的档案、志书、契约、报刊或对方材料。该项锚定的是后续追踪问题，影响范围须按地区与群体分别判断。来源保留：战后、改革后或条约后的记忆、地方记录和官方总结常具有事后选择性；后续影响须按地点、群体和时间分层，不作整体化推断。
\eventanchor{EQ-1789-TIBET-NEPAL-CAMPAIGN-PRELUDE}{「清军越过喜马拉雅方向追击廓…}乾隆五十四年（1789年），「清军越过喜马拉雅方向追击廓尔喀，边境协商与军事行动并行」之前的条件、信息和争议已通过中央与地方材料显现，构成该事件可追踪的前史节点。核心取证：《清高宗实录》乾隆五十四年相关条；军机处档、战事方略及地方奏报；同时期地方档案、地方志、日记或对方材料应按具体地区补检。该项锚定的是前史条件与信息背景，不把条件直接写成唯一因果。来源保留：官修编年和地方材料的记录密度与立场并不相同；本行不将条件、传闻、呈报或后见解释直接等同为确定因果。
\eventanchor{EQ-1789-TIBET-NEPAL-CAMPAIGN-DECISION}{围绕「清军越过喜马拉雅方向追…}乾隆五十四年（1789年），围绕「清军越过喜马拉雅方向追击廓尔喀，边境协商与军事行动并行」的命令、调度、照会或呈报形成独立的中央—地方行动文书链。核心取证：《清高宗实录》乾隆五十四年相关条；军机处档、战事方略及地方奏报。该项锚定的是命令或决策环节，命令抵达和结果仍须另外核验。来源保留：奏折、上谕、部议、军机档或条约可证明行政动作和机构立场，不自动证明所有地区、群体或对手按其意图行动。
\eventanchor{EQ-1789-TIBET-NEPAL-CAMPAIGN-AFTERMATH}{「清军越过喜马拉雅方向追击廓…}乾隆五十四年（1789年），「清军越过喜马拉雅方向追击廓尔喀，边境协商与军事行动并行」引发的善后、执行或地方回应构成可由不同材料追踪的后续节点。核心取证：《清高宗实录》乾隆五十四年相关条；军机处档、战事方略及地方奏报；相关地区的档案、志书、契约、报刊或对方材料。该项锚定的是后续追踪问题，影响范围须按地区与群体分别判断。来源保留：战后、改革后或条约后的记忆、地方记录和官方总结常具有事后选择性；后续影响须按地点、群体和时间分层，不作整体化推断。
\subsection{1790—1799}
\eventanchor{EQ-1790-IMPERIAL-REGULATIONS-TIBET-DOCUMENTS}{清廷围绕「清廷制定整顿西藏事…}乾隆五十五年（1790年），清廷围绕「清廷制定整顿西藏事务的制度性安排，强化驻藏大臣和金瓶掣签等机制」处理盟旗、驻防、交通或地方呈报，相关边务文书进入中央决策链。核心取证：《清高宗实录》乾隆五十五年相关条；军机处档、理藩院档或相关方略。该项锚定的是可回查的文书动作，不将其直接等同于地方执行。来源保留：此行仅确认相关文书链和可追查的行政动作；不得由其直接推断政策有效、地方服从、民众态度或单一因果。
\eventanchor{EQ-1791-WHITE-LOTUS-PRECURSORS-DOCUMENTS}{围绕「川楚陕基层社会的教门网…}乾隆五十六年（1791年），围绕「川楚陕基层社会的教门网络和地方压力成为白莲教战争的背景」的地方呈报、案件或赈务记录将社会反应带入官府文书链。核心取证：《清高宗实录》乾隆五十六年相关条；地方志、督抚奏折及地方档案。该项锚定的是可回查的文书动作，不将其直接等同于地方执行。来源保留：此行仅确认相关文书链和可追查的行政动作；不得由其直接推断政策有效、地方服从、民众态度或单一因果。
\eventanchor{EQ-1792-GURKHA-SETTLEMENT-DOCUMENTS}{围绕「清廷与廓尔喀战争结束后…}乾隆五十七年（1792年），围绕「清廷与廓尔喀战争结束后处理使节、边防和西藏军政善后」的照会、条约文本、换约或地方执行报告分别进入外交文书链。核心取证：《清高宗实录》乾隆五十七年相关条；《筹办夷务始末》相关年卷与中外照会、条约文本。该项锚定的是可回查的文书动作，不将其直接等同于地方执行。来源保留：此行仅确认相关文书链和可追查的行政动作；不得由其直接推断政策有效、地方服从、民众态度或单一因果。
\eventanchor{EQ-1793-MACARTNEY-MISSION-DOCUMENTS}{围绕「马戛尔尼使团抵达清廷，…}乾隆五十八年（1793年），围绕「马戛尔尼使团抵达清廷，围绕贸易、礼仪和外交关系交涉」的照会、条约文本、换约或地方执行报告分别进入外交文书链。核心取证：《清高宗实录》乾隆五十八年相关条；《筹办夷务始末》相关年卷与中外照会、条约文本。该项锚定的是可回查的文书动作，不将其直接等同于地方执行。来源保留：此行仅确认相关文书链和可追查的行政动作；不得由其直接推断政策有效、地方服从、民众态度或单一因果。
\eventanchor{EQ-1794-WHITE-LOTUS-OUTBREAK-DOCUMENTS}{清廷围绕「白莲教相关起事在川…}乾隆五十九年（1794年），清廷围绕「白莲教相关起事在川楚陕地区扩展，长期战争开始」调遣军队、筹措饷械并处理战报，中央与地方的军政文书形成连续记录。核心取证：《清高宗实录》乾隆五十九年相关条；军机处档、战事方略及地方奏报。该项锚定的是可回查的文书动作，不将其直接等同于地方执行。来源保留：此行仅确认相关文书链和可追查的行政动作；不得由其直接推断政策有效、地方服从、民众态度或单一因果。
\eventanchor{EQ-1795-QIANLONG-ABDICATION-DOCUMENTS}{围绕「乾隆帝禅位于嘉庆帝，自…}乾隆六十年（1795年），围绕「乾隆帝禅位于嘉庆帝，自为太上皇，仍保留重要政治影响」的上谕、奏折与部议在中央—地方行政链中流转。核心取证：《清高宗实录》乾隆六十年相关条；宫中朱批奏折、内阁大库题本与《清史稿》相关志传。该项锚定的是可回查的文书动作，不将其直接等同于地方执行。来源保留：此行仅确认相关文书链和可追查的行政动作；不得由其直接推断政策有效、地方服从、民众态度或单一因果。
\eventanchor{EQ-1790-IMPERIAL-REGULATIONS-TIBET-PRELUDE}{「清廷制定整顿西藏事务的制度…}乾隆五十五年（1790年），「清廷制定整顿西藏事务的制度性安排，强化驻藏大臣和金瓶掣签等机制」之前的条件、信息和争议已通过中央与地方材料显现，构成该事件可追踪的前史节点。核心取证：《清高宗实录》乾隆五十五年相关条；军机处档、理藩院档或相关方略；同时期地方档案、地方志、日记或对方材料应按具体地区补检。该项锚定的是前史条件与信息背景，不把条件直接写成唯一因果。来源保留：官修编年和地方材料的记录密度与立场并不相同；本行不将条件、传闻、呈报或后见解释直接等同为确定因果。
\eventanchor{EQ-1790-IMPERIAL-REGULATIONS-TIBET-DECISION}{围绕「清廷制定整顿西藏事务的…}乾隆五十五年（1790年），围绕「清廷制定整顿西藏事务的制度性安排，强化驻藏大臣和金瓶掣签等机制」的命令、调度、照会或呈报形成独立的中央—地方行动文书链。核心取证：《清高宗实录》乾隆五十五年相关条；军机处档、理藩院档或相关方略。该项锚定的是命令或决策环节，命令抵达和结果仍须另外核验。来源保留：奏折、上谕、部议、军机档或条约可证明行政动作和机构立场，不自动证明所有地区、群体或对手按其意图行动。
\eventanchor{EQ-1790-IMPERIAL-REGULATIONS-TIBET-AFTERMATH}{「清廷制定整顿西藏事务的制度…}乾隆五十五年（1790年），「清廷制定整顿西藏事务的制度性安排，强化驻藏大臣和金瓶掣签等机制」引发的善后、执行或地方回应构成可由不同材料追踪的后续节点。核心取证：《清高宗实录》乾隆五十五年相关条；军机处档、理藩院档或相关方略；相关地区的档案、志书、契约、报刊或对方材料。该项锚定的是后续追踪问题，影响范围须按地区与群体分别判断。来源保留：战后、改革后或条约后的记忆、地方记录和官方总结常具有事后选择性；后续影响须按地点、群体和时间分层，不作整体化推断。
\eventanchor{EQ-1791-WHITE-LOTUS-PRECURSORS-PRELUDE}{「川楚陕基层社会的教门网络和…}乾隆五十六年（1791年），「川楚陕基层社会的教门网络和地方压力成为白莲教战争的背景」之前的条件、信息和争议已通过中央与地方材料显现，构成该事件可追踪的前史节点。核心取证：《清高宗实录》乾隆五十六年相关条；地方志、督抚奏折及地方档案；同时期地方档案、地方志、日记或对方材料应按具体地区补检。该项锚定的是前史条件与信息背景，不把条件直接写成唯一因果。来源保留：官修编年和地方材料的记录密度与立场并不相同；本行不将条件、传闻、呈报或后见解释直接等同为确定因果。
\eventanchor{EQ-1791-WHITE-LOTUS-PRECURSORS-DECISION}{围绕「川楚陕基层社会的教门网…}乾隆五十六年（1791年），围绕「川楚陕基层社会的教门网络和地方压力成为白莲教战争的背景」的命令、调度、照会或呈报形成独立的中央—地方行动文书链。核心取证：《清高宗实录》乾隆五十六年相关条；地方志、督抚奏折及地方档案。该项锚定的是命令或决策环节，命令抵达和结果仍须另外核验。来源保留：奏折、上谕、部议、军机档或条约可证明行政动作和机构立场，不自动证明所有地区、群体或对手按其意图行动。
\eventanchor{EQ-1791-WHITE-LOTUS-PRECURSORS-AFTERMATH}{「川楚陕基层社会的教门网络和…}乾隆五十六年（1791年），「川楚陕基层社会的教门网络和地方压力成为白莲教战争的背景」引发的善后、执行或地方回应构成可由不同材料追踪的后续节点。核心取证：《清高宗实录》乾隆五十六年相关条；地方志、督抚奏折及地方档案；相关地区的档案、志书、契约、报刊或对方材料。该项锚定的是后续追踪问题，影响范围须按地区与群体分别判断。来源保留：战后、改革后或条约后的记忆、地方记录和官方总结常具有事后选择性；后续影响须按地点、群体和时间分层，不作整体化推断。
\eventanchor{EQ-1793-MACARTNEY-MISSION-PRELUDE}{「马戛尔尼使团抵达清廷，围绕…}乾隆五十八年（1793年），「马戛尔尼使团抵达清廷，围绕贸易、礼仪和外交关系交涉」之前的条件、信息和争议已通过中央与地方材料显现，构成该事件可追踪的前史节点。核心取证：《清高宗实录》乾隆五十八年相关条；《筹办夷务始末》相关年卷与中外照会、条约文本；同时期地方档案、地方志、日记或对方材料应按具体地区补检。该项锚定的是前史条件与信息背景，不把条件直接写成唯一因果。来源保留：官修编年和地方材料的记录密度与立场并不相同；本行不将条件、传闻、呈报或后见解释直接等同为确定因果。
\eventanchor{EQ-1793-MACARTNEY-MISSION-DECISION}{围绕「马戛尔尼使团抵达清廷，…}乾隆五十八年（1793年），围绕「马戛尔尼使团抵达清廷，围绕贸易、礼仪和外交关系交涉」的命令、调度、照会或呈报形成独立的中央—地方行动文书链。核心取证：《清高宗实录》乾隆五十八年相关条；《筹办夷务始末》相关年卷与中外照会、条约文本。该项锚定的是命令或决策环节，命令抵达和结果仍须另外核验。来源保留：奏折、上谕、部议、军机档或条约可证明行政动作和机构立场，不自动证明所有地区、群体或对手按其意图行动。
\eventanchor{EQ-1793-MACARTNEY-MISSION-AFTERMATH}{「马戛尔尼使团抵达清廷，围绕…}乾隆五十八年（1793年），「马戛尔尼使团抵达清廷，围绕贸易、礼仪和外交关系交涉」引发的善后、执行或地方回应构成可由不同材料追踪的后续节点。核心取证：《清高宗实录》乾隆五十八年相关条；《筹办夷务始末》相关年卷与中外照会、条约文本；相关地区的档案、志书、契约、报刊或对方材料。该项锚定的是后续追踪问题，影响范围须按地区与群体分别判断。来源保留：战后、改革后或条约后的记忆、地方记录和官方总结常具有事后选择性；后续影响须按地点、群体和时间分层，不作整体化推断。
\eventanchor{EQ-1794-WHITE-LOTUS-OUTBREAK-PRELUDE}{「白莲教相关起事在川楚陕地区…}乾隆五十九年（1794年），「白莲教相关起事在川楚陕地区扩展，长期战争开始」之前的条件、信息和争议已通过中央与地方材料显现，构成该事件可追踪的前史节点。核心取证：《清高宗实录》乾隆五十九年相关条；军机处档、战事方略及地方奏报；同时期地方档案、地方志、日记或对方材料应按具体地区补检。该项锚定的是前史条件与信息背景，不把条件直接写成唯一因果。来源保留：官修编年和地方材料的记录密度与立场并不相同；本行不将条件、传闻、呈报或后见解释直接等同为确定因果。
\eventanchor{EQ-1794-WHITE-LOTUS-OUTBREAK-DECISION}{围绕「白莲教相关起事在川楚陕…}乾隆五十九年（1794年），围绕「白莲教相关起事在川楚陕地区扩展，长期战争开始」的命令、调度、照会或呈报形成独立的中央—地方行动文书链。核心取证：《清高宗实录》乾隆五十九年相关条；军机处档、战事方略及地方奏报。该项锚定的是命令或决策环节，命令抵达和结果仍须另外核验。来源保留：奏折、上谕、部议、军机档或条约可证明行政动作和机构立场，不自动证明所有地区、群体或对手按其意图行动。
\eventanchor{EQ-1794-WHITE-LOTUS-OUTBREAK-AFTERMATH}{「白莲教相关起事在川楚陕地区…}乾隆五十九年（1794年），「白莲教相关起事在川楚陕地区扩展，长期战争开始」引发的善后、执行或地方回应构成可由不同材料追踪的后续节点。核心取证：《清高宗实录》乾隆五十九年相关条；军机处档、战事方略及地方奏报；相关地区的档案、志书、契约、报刊或对方材料。该项锚定的是后续追踪问题，影响范围须按地区与群体分别判断。来源保留：战后、改革后或条约后的记忆、地方记录和官方总结常具有事后选择性；后续影响须按地点、群体和时间分层，不作整体化推断。
\eventanchor{EQ-1795-QIANLONG-ABDICATION-PRELUDE}{「乾隆帝禅位于嘉庆帝，自为太…}乾隆六十年（1795年），「乾隆帝禅位于嘉庆帝，自为太上皇，仍保留重要政治影响」之前的条件、信息和争议已通过中央与地方材料显现，构成该事件可追踪的前史节点。核心取证：《清高宗实录》乾隆六十年相关条；宫中朱批奏折、内阁大库题本与《清史稿》相关志传；同时期地方档案、地方志、日记或对方材料应按具体地区补检。该项锚定的是前史条件与信息背景，不把条件直接写成唯一因果。来源保留：官修编年和地方材料的记录密度与立场并不相同；本行不将条件、传闻、呈报或后见解释直接等同为确定因果。
\eventanchor{EQ-1795-QIANLONG-ABDICATION-DECISION}{围绕「乾隆帝禅位于嘉庆帝，自…}乾隆六十年（1795年），围绕「乾隆帝禅位于嘉庆帝，自为太上皇，仍保留重要政治影响」的命令、调度、照会或呈报形成独立的中央—地方行动文书链。核心取证：《清高宗实录》乾隆六十年相关条；宫中朱批奏折、内阁大库题本与《清史稿》相关志传。该项锚定的是命令或决策环节，命令抵达和结果仍须另外核验。来源保留：奏折、上谕、部议、军机档或条约可证明行政动作和机构立场，不自动证明所有地区、群体或对手按其意图行动。
\eventanchor{EQ-1795-QIANLONG-ABDICATION-AFTERMATH}{「乾隆帝禅位于嘉庆帝，自为太…}乾隆六十年（1795年），「乾隆帝禅位于嘉庆帝，自为太上皇，仍保留重要政治影响」引发的善后、执行或地方回应构成可由不同材料追踪的后续节点。核心取证：《清高宗实录》乾隆六十年相关条；宫中朱批奏折、内阁大库题本与《清史稿》相关志传；相关地区的档案、志书、契约、报刊或对方材料。该项锚定的是后续追踪问题，影响范围须按地区与群体分别判断。来源保留：战后、改革后或条约后的记忆、地方记录和官方总结常具有事后选择性；后续影响须按地点、群体和时间分层，不作整体化推断。
